\documentclass[
    12pt
    ,oneside
    ,a4paper
    ,chapter=TITLE
    ,section=TITLE
    ,sumario=abnt-6027-2012]{abntex2}

% Carrega o pacote da abnt com as adapatções para normas da Tuiuti
\usepackage{packages/abnt-UTP}

% =================================================================================================
% -------------------------------------------------------------------------------------------------
% -                                                                                               -
% - INFORMAÇÕES GERAIS                                                                            -
% -                                                                                               -
% -------------------------------------------------------------------------------------------------
% =================================================================================================

% Informações para CAPA e FOLHA DE ROSTO
\titulo{Desenvolvimento de um Software com Inteligência Artificial Explicável (XAI) para Apoio ao Diagnóstico Preditivo de Câncer de Mama utilizando SHAP e UMAP}

\autor{Gabriel Ast dos Santos}

\orientador{Prof. Rodrigo Ramos Alves}

\preambulo{Trabalho de Conclusão de Curso apresentado ao curso de Bacharelado em Ciência da
Computação da Faculdade de Ciências Exatas e de Tecnologia da Universidade Tuiuti do Paraná,
como requisito à obtenção ao grau de Bacharel.}

\instituicao{Universidade Tuiuti do Paraná}
\local{Curitiba}
\data{\the\year}

% =================================================================================================
% -------------------------------------------------------------------------------------------------
% -                                                                                               -
% - DOCUMENTO                                                                                     -
% -                                                                                               -
% -------------------------------------------------------------------------------------------------
% =================================================================================================
\begin{document}

% =================================================================================================
% -------------------------------------------------------------------------------------------------
% - ELEMENTOS PRÉ-TEXTUAIS                                                                        -
% -------------------------------------------------------------------------------------------------
% =================================================================================================

% -------------------------------------------------------------------------------------------------
% Cria a capa e a folha de rosto
\imprimircapa
\imprimirfolhaderosto

% % Resumo do trabalho
\begin{resumo}
    O câncer de mama é a neoplasia mais incidente entre mulheres no Brasil, representando 30\% dos novos diagnósticos femininos, o que torna o diagnóstico precoce um fator determinante para a eficácia do tratamento. Este trabalho propõe o desenvolvimento de um software integrado a um modelo preditivo para a classificação de tumores mamários como benignos ou malignos, utilizando a base de dados pública Breast Cancer Wisconsin. O objetivo central é mitigar a opacidade característica dos modelos de caixa preta, fornecendo diagnósticos que unam alta precisão analítica a justificativas clínicas transparentes e interpretáveis por profissionais de saúde. A metodologia contempla o treinamento e a comparação de algoritmos clássicos de aprendizado de máquina supervisionado, incluindo Regressão Logística, Máquina de Vetores de Suporte, K-Nearest Neighbors, Random Forest e XGBoost, avaliados por métricas de sensibilidade, precisão, F1-score e área sob a curva de precisão-revocação. A explicabilidade das decisões é garantida pela integração da biblioteca SHapley Additive exPlanations (SHAP), que quantifica a contribuição individual de cada variável celular morfológica no risco de malignidade apontado pelo modelo, por meio de visualizações locais e globais. Em paralelo, o algoritmo Uniform Manifold Approximation and Projection (UMAP) é aplicado para reduzir a dimensionalidade dos dados e revelar graficamente a separabilidade natural entre tumores benignos e malignos, facilitando a identificação de agrupamentos e anomalias biológicas. Por fim, o sistema disponibiliza uma interface capaz de receber dados em lote via upload de planilhas e gerar relatórios estruturados com os diagnósticos e suas respectivas fundamentações visuais. Os resultados esperados indicam que a combinação de modelos preditivos de alta performance com técnicas de inteligência artificial explicável constitui uma solução viável e eticamente fundamentada para o apoio ao diagnóstico oncológico assistido por computador.
    
    \palavraschave{Aprendizado de máquina; Base de dados Wisconsin; Câncer de mama; Classificação binária; Inteligência Artificial Explicável; SHAP; UMAP.}
\end{resumo}


% % Listas de ilustrações, tabelas, siglas, etc.
% % - Inclua apenas as listas dos elementos que serão utilizados no trabalho.
% % - A lista de símbolos e siglas devem ser editadas manualmente 
\listadefiguras
% \listadegraficos
\listadetabelas
% \listadequadros
% \listadecodigos
% \listadealgoritmos


% Lista de siglas
% - As siglas deve ser organizadas em ORDEM ALFABÉTICA
% - Inclua apenas as siglas que aparecem no trabalho
\begin{siglas}
    \item[AUC] \textit{Area Under the Curve} (Área Sob a Curva)
    \item[DNA] Ácido desoxirribonucleico
    \item[DL] \textit{Deep Learning}
    \item[IA] Inteligência Artificial
    \item[INCA] Instituto Nacional de Câncer
    \item[KNN] \textit{K-Nearest Neighbors} (K-Vizinhos Mais Próximos)
    \item[ML] \textit{Machine Learning}
    \item[PAAF] Punção Aspirativa por Agulha Fina
    \item[PAHO] \textit{Pan American Health Organization}
    \item[RBF] \textit{Radial Basis Function} (Função de Base Radial) 
    \item[ROC] \textit{Receiver Operating Characteristic} (Curva de Característica de Operação do Receptor)
    \item[SHAP] \textit{SHapley Additive exPlanations}
    \item[SVM] \textit{Support Vector Machine} (Máquina de Vetores de Suporte)
    \item[UMAP] \textit{Uniform Manifold Approximation and Projection}
    \item[XAI] \textit{eXplainable AI} (Inteligência Artificial Explicável)
\end{siglas}



% Lista de símbolos
\begin{simbolos}
    \item[$ \phi_i $]   Valor de Shapley para a característica isolada $i$
    \item[$ \Sigma $]   Somatório
    \item[$ N $]        Conjunto total de características celulares disponíveis
    \item[$ S $]        Subconjunto específico de características
    \item[$ i $]        Característica celular isolada em análise
    \item[$ \subseteq $] Subconjunto de
    \item[$ \setminus $] Diferença de conjuntos (exclusão)
    \item[$ | \cdot | $] Cardinalidade (quantidade de elementos de um conjunto)
    \item[$ ! $]        Operador fatorial
    \item[$ \cup $]     União de conjuntos
    \item[$ v(\cdot) $] Função de valor (predição de risco do modelo para um determinado subconjunto)
\end{simbolos}

% % Criação do sumário do trabalho
\sumario

% =================================================================================================
% -------------------------------------------------------------------------------------------------
% - ELEMENTOS TEXTUAIS                                                                            -
% -------------------------------------------------------------------------------------------------
% =================================================================================================

% Reinicia a numeração das páginas
\textual

% -------------------------------------------------------------------------------------------------
% -------------------------------------------------------------------------------------------------
% 1. INTRODUÇÃO
% -------------------------------------------------------------------------------------------------
% -------------------------------------------------------------------------------------------------
\chapter{Introdução}
\label{cap:introducao}

De acordo com a Organização Pan-Americana da Saúde \cite{Paho-cancer}, as patologias oncológicas impõem uma enorme carga de adoecimento nas Américas, sendo superadas unicamente pelas doenças cardiovasculares. A instituição aponta que, durante o ano de 2022, a região contabilizou mais de 4,2 milhões de novos casos, com a expectativa de que esse índice sofra um acréscimo de 60\% até 2045, chegando à marca de 6,7 milhões. Como uma das causas primárias de mortalidade no continente, a doença resultou em 1,4 milhão de perdas fatais em 2022, afetando pessoas com 69 anos ou menos em 45\% das ocorrências. Contudo, o órgão ressalta que muitos desses quadros possuem elevadas taxas de cura quando há o rastreio precoce somado ao tratamento clínico assertivo.

No cenário nacional, as projeções apontam para um panorama de alerta, com a estimativa de 781 mil novos casos de câncer por ano no Brasil ao longo do triênio de 2026 a 2028. Dentro dessa estatística, o câncer de mama consolida-se de forma alarmante como a neoplasia mais incidente entre o público feminino, representando 30\% de todos os novos diagnósticos nas mulheres. Esses números expressivos refletem não apenas o envelhecimento populacional, mas também evidenciam as profundas desigualdades regionais e os desafios persistentes que as pacientes enfrentam no acesso à prevenção, ao rastreamento, ao diagnóstico precoce e ao tratamento oportuno no sistema de saúde \cite{Inca-781}.

O uso de tecnologias computacionais para enfrentar essa carga de adoecimento está em plena expansão, impulsionado pela chamada revolução do \textit{Big Data} na saúde. Instituições de referência, como o próprio Instituto Nacional de Câncer \cite{Inca-informatica}, destacam que a bioinformática e a biologia computacional são cruciais para analisar dados em larga escala e identificar assinaturas do câncer. Ao processar grandes volumes de dados morfológicos e genéticos, abrem-se caminhos para uma medicina cada vez mais direcionada e personalizada. Nesse contexto, o uso de algoritmos de \textit{Machine Learning} para o reconhecimento de padrões e classificação de tumores já se consolida como uma linha de pesquisa fundamental para a oncologia moderna, capaz de extrair marcadores de interesse clínico de forma automatizada.

Apesar do notável desempenho preditivo, a adoção clínica dessas tecnologias esbarra no chamado "problema da caixa preta". Algoritmos avançados fornecem predições altamente acuradas, mas não revelam o raciocínio matemático interno que levou àquela conclusão. Na prática médica, entregar um diagnóstico de malignidade automatizado sem uma fundamentação biológica rastreável é insuficiente e eticamente questionável. O profissional de saúde necessita compreender quais fatores morfológicos e celulares pesaram na decisão para validar a predição, confiar no modelo e transmitir o diagnóstico com segurança ao paciente.

É nesse cenário que a Inteligência Artificial Explicável (XAI - \textit{eXplainable AI}) se faz estritamente necessária. A aplicação de técnicas de interpretabilidade, como o \textit{SHapley Additive exPlanations} (SHAP), permite "abrir" essa caixa preta ao quantificar o impacto exato e a contribuição de cada variável clínica no risco de malignidade. Em paralelo, a utilização de algoritmos de redução de dimensionalidade, como o \textit{Uniform Manifold Approximation and Projection for Dimension Reduction} (UMAP), facilita a análise visual da separabilidade natural dos dados, identificando agrupamentos (\textit{clusters}) e anomalias de forma gráfica. Dessa forma, o sistema computacional deixa de ser apenas um classificador opaco e passa a atuar como uma ferramenta de apoio transparente, unindo a precisão analítica à explicabilidade exigida pelo campo oncológico.

Diante disso, surge o seguinte problema de pesquisa: como integrar técnicas de explicabilidade e redução de dimensionalidade em um modelo preditivo de câncer de mama, de modo a fornecer diagnósticos que sejam, simultaneamente, acurados e interpretáveis por profissionais de saúde?

O objetivo deste trabalho é desenvolver um software integrado a um modelo preditivo para a classificação de tumores mamários (benignos ou malignos), utilizando algoritmos de Machine Learning e a base de dados pública Breast Cancer Wisconsin. O intuito principal é mitigar a opacidade característica dos modelos de "caixa preta", fornecendo aos profissionais de saúde diagnósticos preditivos que unam alta precisão analítica a justificativas clínicas transparentes.

Para alcançar esse propósito, o projeto engloba o treinamento e a seleção de classificadores, bem como a aplicação do UMAP para a análise visual da separabilidade natural dos dados, facilitando a identificação de agrupamentos (clusters) e anomalias (outliers) biológicas. Em paralelo, o sistema integrará a biblioteca SHAP para garantir a explicabilidade das decisões, quantificando o impacto de cada variável celular no risco de malignidade apontado pelo modelo. Por fim, visando a viabilidade de uso no cotidiano médico, o trabalho contempla a criação de uma interface capaz de receber dados em lote via upload de planilhas, automatizando a geração de relatórios de saída estruturados (como PDF, Excel ou CSV) com os resumos dos diagnósticos e suas respectivas fundamentações visuais.


% -------------------------------------------------------------------------------------------------
% -------------------------------------------------------------------------------------------------
% 2. FUNDAMENTAÇÃO TEÓRICA
% -------------------------------------------------------------------------------------------------
% -------------------------------------------------------------------------------------------------
\chapter{Fundamentação}
\label{cap:fundamentacao}

Este capítulo destina-se à fundamentação teórica que sustenta a pesquisa, conectando conceitos fundamentais ao problema proposto. O objetivo é estabelecer uma base sólida, com base na literatura científica consolidada, que conduza e valide o estudo. Nas seções a seguir, serão explorados os aspectos clínicos da doença, a evolução tecnológica na análise de dados em saúde e as abordagens computacionais escolhidas para o desenvolvimento do modelo preditivo.

\section{O Câncer de Mama e a Análise Morfológica}

O termo câncer abrange um complexo conjunto de mais de cem doenças malignas que têm como denominador comum o crescimento desordenado e incontrolável das células. A base biológica dessa patologia reside na perda estrutural e funcional dos mecanismos de regulação celular. Esse processo de formação tumoral, cientificamente denominado carcinogênese ou oncogênese, é engatilhado quando o ácido desoxirribonucleico (DNA), a "memória química"\ que guarda as instruções das atividades celulares, sofre alterações críticas. Em condições normais, genes reguladores específicos chamados proto-oncogenes permanecem inativos; contudo, sob a ação de fatores cancerígenos cumulativos, eles sofrem mutações e são ativados, convertendo-se em oncogenes. Essa ativação desencadeia um processo lento e progressivo composto por três estágios: a iniciação, onde a célula sofre a alteração genética primária de forma silenciosa, a promoção, fase em que agentes externos ou internos transformam gradual e lentamente essa célula iniciada em uma estrutura maligna, e por fim, a progressão, caracterizada pela multiplicação descontrolada e irreversível das células alteradas, culminando na consolidação clínica do tumor e em seu potencial de invasão \cite{Inca-como-surge-cancer}.

O tipo mais comum desta doença é classificado como carcinoma, uma nomeação atribuída aos tumores malignos que se originam nas células epiteliais, que são os tecidos responsáveis por revestir as superfícies internas e externas do corpo humano. Dependendo do tipo específico de célula epitelial afetada, o carcinoma recebe denominações distintas. No contexto anatômico do câncer de mama, a manifestação predominante é o adenocarcinoma, um subtipo que se desenvolve especificamente nas células epiteliais dos tecidos glandulares, cuja função biológica é produzir fluidos e muco \cite{Nci-what-is-cancer}.

No entanto, o desenvolvimento desse adenocarcinoma não ocorre de maneira instantânea, estruturando-se por meio de uma progressão morfológica gradual e microscopicamente rastreável. Inicialmente, o tecido mamário pode apresentar uma hiperplasia, caracterizada por um acúmulo de células que se multiplicam em um ritmo acelerado, mas que ainda preservam rigorosamente sua aparência e arquitetura originais. Se a anomalia celular progredir para o estágio de displasia, além do acúmulo excessivo, as células perdem sua organização e começam a apresentar deformidades morfológicas evidentes. O ápice dessa fase pré-invasiva é denominado carcinoma \textit{in situ} (frequentemente referido na literatura clínica como estágio 0). Neste limiar, as células já exibem características visualmente malignas em seus núcleos, mas ainda se encontram confinadas, pois não romperam as barreiras do tecido glandular original para invadir as áreas vizinhas \cite{Nci-what-is-cancer}. A diferença fundamental entre a biologia de um tecido saudável e o desenvolvimento anômalo gerado por essa cascata de mutações pode ser observada na Figura \ref{fig:divisao_celular}.

\begin{figure}[htb]
    \centering
    \caption{Contraste morfológico entre células normais e neoplásicas.}
    \label{fig:divisao_celular}
    \includegraphics[width=0.8\textwidth]{oqueecancercelulas.jpg} 
    \fonte{\citeauthoronline{Inca-O-que-e-cancer}, \citeyear{Inca-O-que-e-cancer}, Disponível em: \url{https://www.gov.br/inca/pt-br/assuntos/cancer/o-que-e-cancer}}
\end{figure}

Na ciência da computação aplicada à saúde, o dado mais valioso desta anomalia não é apenas a mutação genética em si, mas sim a profunda alteração morfológica que ela desencadeia. A pressão para sustentar uma replicação em massa e desordenada leva à deformação severa do núcleo de uma célula maligna. Como resultado, quando analisadas ao microscópio, estas células apresentam variações perceptíveis e quantificáveis na área, no raio, na textura e na concavidade, em comparação com uma célula saudável. São justamente essas assimetrias físicas que se tornam o insumo fundamental para alimentar e treinar algoritmos de aprendizagem de máquina. Estes algoritmos apoiam o diagnóstico preditivo.

Um dos marcos pioneiros na convergência entre a oncologia clínica e a ciência da computação, voltado especificamente para a quantificação dessas assimetrias celulares, é a base de dados \textit{Breast Cancer Wisconsin (Diagnostic)}. No início dos anos 90, pesquisadores da Universidade de Wisconsin criaram esta base para reduzir a subjetividade e o cansaço que os humanos sentem quando analisam coisas ao microscópio \cite{Street1993}. O mais importante neste trabalho foi a decisão de não criar um armazém de imagens de doença com muitos pixels, que precisariam de muita capacidade de processamento, mas sim uma base de dados organizada e que não precisasse de muita capacidade de computação.

Para alcançar esse resultado, os pesquisadores utilizaram amostras coletadas via Punção Aspirativa por Agulha Fina (PAAF) e aplicaram algoritmos de processamento de imagem, conhecidos como contornos ativos (\textit{snakes}), para delinear digitalmente as fronteiras dos núcleos celulares. A partir deste tratamento computacional, foram extraídas matematicamente dez características morfológicas essenciais de cada célula, incluindo o raio, a textura, o perímetro, a área e a suavidade. Ao calcular a média, o erro padrão e o valor máximo para cada uma dessas dez medidas, o sistema gerou as 30 variáveis numéricas que compõem o conjunto de dados. Esta otimização consolidou a base de Wisconsin como um recurso de referência na literatura, permitindo que os algoritmos preditivos processem informações oncológicas complexas de forma rápida e eficiente.

\section{Inteligência Artificial e Machine Learning: Definições e Hierarquia}

A Inteligência Artificial (IA) é uma área da ciência da computação que começou no verão de 1956, na conferência de Dartmouth, onde o termo foi criado. Desde então, a área evoluiu e ramificou-se em diversas vertentes de pesquisa. Segundo \citeonline{Russell2010}, as definições históricas e contemporâneas de IA podem ser organizadas em quatro categorias principais, divididas entre abordagens que buscam simular o comportamento ou o pensamento humano, e abordagens focadas puramente no conceito de racionalidade. No contexto computacional moderno, a abordagem mais adotada e cientificamente rigorosa é a dos "agentes racionais", que são sistemas que operam de forma autônoma e agem para alcançar o melhor resultado possível diante de um problema, dado o conhecimento e as percepções que possuem do ambiente.

Para que um sistema alcance essa racionalidade em ambientes complexos, a programação de regras lógicas estáticas (\textit{if-then}) torna-se inviável. É para solucionar esta limitação que surge o \textit{Machine Learning} (ML), o principal subcampo da Inteligência Artificial. Em vez de serem explicitamente programados para executar uma tarefa passo a passo, os algoritmos de ML utilizam métodos estatísticos para extrair padrões diretamente dos dados. A definição formal mais consolidada na literatura é proposta por \citeonline{Mitchell1997}, que estabelece que um programa de computador aprende a partir da experiência $E$, em relação a uma classe de tarefas $T$ e uma medida de desempenho $P$, se o seu desempenho em $T$, quantificado por $P$, melhora com a experiência $E$.

Para compreender essa definição de forma mais prática, pode-se estabelecer um paralelo com o aprendizado humano em jogos. A tarefa $T$ representa o objetivo, como identificar tumores; a experiência $E$ consiste nos dados históricos e exemplos fornecidos ao sistema (como a base de dados Wisconsin), e o desempenho $P$ atua como a métrica de sucesso, indicando quantos diagnósticos a máquina acertou. O algoritmo é considerado capaz de "aprender"\ quando, ao ser exposto a mais dados (maior $E$), ele aumenta sua taxa de acertos (maior $P$) ao realizar o diagnóstico (tarefa $T$).

A hierarquia da inteligência computacional estende-se ainda a camadas mais especializadas, como o \textit{Deep Learning} (DL), ou Aprendizado Profundo. Essa vertente é um subconjunto do ML que utiliza arquiteturas de redes neurais artificiais com múltiplas camadas para processar volumes massivos de dados de altíssima complexidade dimensional. Essa relação hierárquica e de pertencimento entre os campos pode ser visualizada por meio de um diagrama de conjuntos, conforme ilustrado na Figura \ref{fig:hierarquia_ia}. 

\begin{figure}[htb]
    \centering
    \caption{Hierarquia e relação de inclusão entre IA, ML e DL.}
    \label{fig:hierarquia_ia}
    \includegraphics[width=0.6\textwidth]{venn_ia_ml_dp.jpg} 
    \fonte{\citeauthoronline{RedHat-IAML}, \citeyear{RedHat-IAML}, Disponível em: \url{https://www.redhat.com/pt-br/blog/what-aiml-and-why-does-it-matter-your-business}}
\end{figure}

Essa estrutura demonstra que todo Aprendizado Profundo é uma forma de Aprendizado de Máquina, e todo Aprendizado de Máquina é uma aplicação de Inteligência Artificial. Para o escopo do diagnóstico preditivo oncológico proposto neste estudo, a construção do agente racional será fundamentada especificamente nas técnicas clássicas de \textit{Machine Learning}, cujos paradigmas e mecânicas de avaliação serão detalhados a seguir.

\section{Machine Learning: Conceitos e Paradigmas}
% Aprendizado Supervisionado. Classificação Binária. O Ciclo de Vida: Coleta, pré-processamento, treinamento e avaliação. Desafios: Overfitting, underfitting, viés e variância. Importância da validação cruzada e do conjunto de teste. Métricas de avaliação: acurácia, precisão, recall, F1-score, AUC-ROC. Algoritmos comuns: Regressão Logística, SVM, Random Forest, XGBoost, Redes Neurais. 

O desenvolvimento de um modelo de IA não segue uma receita universal, a abordagem escolhida depende fundamentalmente da natureza dos dados disponíveis e do objetivo do problema. Na literatura clássica de \textit{Machine Learning}, os algoritmos são categorizados em três grandes paradigmas de aprendizado: Não Supervisionado, Aprendizado por Reforço e Aprendizado Supervisionado. Como o escopo deste trabalho visa prever um diagnóstico clínico a partir de dados históricos previamente validados por médicos patologistas, o referencial teórico e metodológico desta pesquisa concentra-se exclusivamente no paradigma do Aprendizado Supervisionado.

\subsection{O Paradigma do Aprendizado Supervisionado e a Classificação Binária}

O Aprendizado Supervisionado, que a literatura classifica como uma tarefa preditiva de aprendizagem indutiva, caracteriza-se sobretudo pela presença de um "gabarito"\ nos dados de treinamento. Neste paradigma, o algoritmo é alimentado com um conjunto de dados históricos contendo tanto as variáveis independentes de entrada, conhecidas como \textit{features} ou atributos, quanto as suas respectivas variáveis dependentes de saída, os rótulos ou \textit{labels}. O objetivo deste método é fazer com que o modelo descubra e mapeie a relação matemática entre as entradas e a saída. Dessa forma, ao aprender as regras ocultas que conectam as características observadas ao resultado final, o sistema torna-se capaz de generalizar esse conhecimento para realizar predições automáticas quando exposto a novos dados não rotulados no futuro \cite{Faceli2011}.

Em termos didáticos, o processo assemelha-se a um professor ensinando um aluno por meio de exemplos. O algoritmo analisa casos em que tanto o problema quanto a resposta correta são fornecidos simultaneamente. Após essa fase de treinamento, o modelo deve ser capaz de generalizar o conhecimento adquirido para prever a resposta correta quando apresentado a novos dados inéditos, que não possuem o rótulo. 

Dentro do Aprendizado Supervisionado, \citeonline{Bishop2006} divide os problemas de acordo com o tipo de resultado que se deseja alcançar. Quando o objetivo do modelo é organizar as informações recebidas em grupos bem definidos, a tarefa é chamada de problema de classificação. Por outro lado, se o sistema tenta prever um valor numérico contínuo, como estimar a expectativa de vida de um paciente ou calcular o rendimento de um processo químico, a tarefa recebe o nome de problema de regressão.

No contexto da oncologia e da base de dados \textit{Breast Cancer Wisconsin}, a tarefa computacional se enquadra como um problema de Classificação Binária. Isso significa que a variável de saída que o modelo deve prever possui apenas dois estados: o tumor é Benigno ou Maligno \cite{Street1993}. Todo o esforço de treinamento e otimização dos algoritmos neste trabalho visa buscar a melhor fronteira de decisão matemática capaz de separar essas duas classes, minimizando ao máximo os erros de categorização, dada a criticidade de um falso diagnóstico médico.

\subsection{O Dilema da Generalização e o Ciclo de Vida do Modelo}

A construção de um classificador médico eficiente vai muito além do simples processamento numérico de dados históricos. O objetivo central de um modelo de aprendizado de máquina não é apenas memorizar os exemplos exatos que foram apresentados durante o treinamento, mas sim desenvolver a capacidade de generalização. Ou seja, o sistema precisa aprender a regra geral para conseguir prever com precisão o diagnóstico de novos pacientes que nunca processou antes. Nesse processo de aprendizado, o algoritmo pode cometer dois erros estatísticos extremos: o \textit{underfitting} e o \textit{overfitting}.

O \textit{underfitting} ocorre quando o modelo matemático é estruturalmente simples demais e não consegue capturar nem mesmo os padrões morfológicos básicos dos tumores. Nesse cenário, o algoritmo falha em extrair o conhecimento necessário e erra classificações fundamentais. No extremo oposto temos o \textit{overfitting}, que é um risco estatístico muito mais comum e perigoso na área da saúde. Nesse caso, o modelo se torna tão complexo que passa a "decorar"\ não apenas as características reais do câncer, mas também as falhas, as exceções e os pequenos ruídos específicos daquela base de dados de treinamento. Consequentemente, o modelo atinge uma precisão artificialmente alta durante seu desenvolvimento, mas fracassa ao avaliar novos pacientes, pois apenas memorizou os dados antigos em vez de assimilar a lógica subjacente à doença \cite{Faceli2011}.

Na estatística, essa linha tênue entre memorizar e generalizar é explicada pelo dilema viés e variância (\textit{bias-variance tradeoff}). O viés representa a limitação estrutural do modelo em enxergar a complexidade dos dados reais, o que causa o \textit{underfitting}. Já a variância mede o quanto o modelo altera sua fronteira de decisão de forma exagerada diante de pequenas mudanças nos dados, o que resulta no \textit{overfitting}. Para garantir que o classificador encontre o equilíbrio estatístico perfeito e seja seguro para uso clínico, é obrigatório adotar um ciclo de vida rigoroso no qual a base de dados original nunca é entregue inteira de uma só vez ao algoritmo. Ela precisa ser metodologicamente dividida em subconjuntos separados \cite{Bishop2006}.

Essa metodologia de divisão geralmente particiona os dados em três grupos: Treinamento, Validação e Teste. O conjunto de treinamento é o ambiente computacional primário, onde o algoritmo ajusta seus pesos e constrói suas regras de classificação. O conjunto de validação atua como um ambiente de aferição intermediária, permitindo que o desenvolvedor calibre os detalhes matemáticos do modelo e perceba rapidamente se o \textit{overfitting} está ocorrendo. Por fim, o conjunto de teste representa o cenário do mundo real. Ele permanece totalmente isolado e oculto do computador durante todo o processo de desenvolvimento, sendo acionado apenas uma única vez na etapa final. O desempenho alcançado neste conjunto de teste isolado é o indicador soberano que comprova a verdadeira eficácia do sistema em acertar diagnósticos na prática médica.

\subsection{Métricas de Avaliação de Desempenho}

% Matriz de Confusão
Após submeter o conjunto de teste ao algoritmo, é fundamental quantificar o seu desempenho clínico. Na classificação binária, a base para qualquer avaliação é a construção de uma Matriz de Confusão. Essa ferramenta organiza os resultados das previsões em quatro categorias possíveis. Os Verdadeiros Positivos (VP) representam os tumores malignos corretamente identificados, enquanto os Verdadeiros Negativos (VN) são os tumores benignos confirmados. Por outro lado, os erros são divididos em Falsos Positivos (FP), que ocorrem quando um tumor benigno gera um alarme falso de malignidade, e Falsos Negativos (FN), que acontecem quando o sistema ignora um tumor maligno e o classifica como inofensivo. Em um cenário oncológico, o falso negativo representa a falha mais crítica do modelo, pois resulta em um paciente doente que é liberado sem receber o tratamento adequado \cite{James2023}.

A Figura~\ref{fig:matriz_confusao} ilustra um exemplo dessa estrutura, onde os acertos do modelo concentram-se na diagonal principal (430 VN e 61 VP), enquanto os erros aparecem nos quadrantes secundários (6 FP e 3 FN).
\begin{figure}[htpb]
    \centering
    \caption{Exemplo Ilustrativo de uma Matriz de Confusão}
    \includegraphics[width=0.55\textwidth]{CM_Exemplo.png}
    \label{fig:matriz_confusao}
    \fonte{O Autor}
\end{figure}



% Acurácia
A métrica mais intuitiva extraída dessa matriz é a acurácia (\textit{accuracy}), que representa a taxa geral de acertos do modelo. Contudo, confiar exclusivamente na acurácia é um erro metodológico perigoso na área da saúde devido ao problema matemático das classes desbalanceadas. Considere uma situação onde a base de dados possui noventa e nove pacientes saudáveis para cada paciente com câncer. Se um algoritmo quebrado simplesmente classificar todas as pessoas do mundo como saudáveis, ele obterá uma acurácia de 99\%. Apesar de parecer um sucesso estatístico absoluto, o modelo seria clinicamente inútil e letal para a pequena parcela de pacientes doentes \cite{Faceli2011}. Matematicamente, a acurácia é calculada pela razão entre as predições corretas e o total de predições, conforme a Equação \ref{eq:acuracia}. Seu resultado é expresso em uma escala de 0 a 1, onde valores próximos a 1 indicam maior acerto global.
\begin{equation}
    \label{eq:acuracia}
    Acur\acute{a}cia = \frac{VP + VN}{VP + VN + FP + FN}
\end{equation}



% Sensibilidade/Recall
Para contornar essa ilusão estatística, o desempenho do modelo deve ser avaliado por métricas que observam os acertos e erros de cada classe de forma isolada. A sensibilidade, também conhecida na literatura como revocação (\textit{recall}), mede a capacidade do algoritmo de identificar corretamente os pacientes que realmente possuem a doença, sendo expressa pela Equação \ref{eq:sensibilidade} \cite{Davis2006}. Em outras palavras, essa métrica quantifica a proporção de casos positivos que são efetivamente recuperados pelo modelo, sendo particularmente relevante em contextos clínicos onde a omissão de diagnósticos pode acarretar consequências graves.

\begin{equation}
    \label{eq:sensibilidade}
    Sensibilidade = \frac{VP}{VP + FN}
\end{equation}



% Precisão
Em paralelo, a precisão (\textit{precision}) avalia a confiabilidade técnica do alerta, ou seja, de todos os tumores que o sistema afirmou serem malignos, qual proporção realmente era, expressada pela Equação \ref{eq:precisao} \cite{Davis2006}. Essa métrica está diretamente relacionada à taxa de alarmes falsos do sistema, refletindo o quanto as predições positivas podem ser consideradas confiáveis na prática. Assim, modelos com alta precisão tendem a evitar intervenções desnecessárias em pacientes saudáveis.

\begin{equation}
    \label{eq:precisao}
    Precis\tilde{a}o = \frac{VP}{VP + FP}
\end{equation}



% F1-Score
Como frequentemente existe uma balança entre essas duas métricas, onde o aumento de uma pode diminuir a outra, utiliza-se a métrica \textit{F1-score}. Esse indicador calcula a média harmônica entre a precisão e a sensibilidade, fornecendo um único número que atesta o equilíbrio e a segurança geral do classificador \cite{Manning2006}. O \textit{F1-score} consolida esses valores através da Equação \ref{eq:f1score}.

\begin{equation}
    \label{eq:f1score}
    F1\text{-}Score = 2 \times \frac{Precis\tilde{a}o \times Sensibilidade}{Precis\tilde{a}o + Sensibilidade}
\end{equation}



% ROC e AUC
A avaliação final do modelo é frequentemente consolidada visualmente pela Curva de Característica de Operação do Receptor (ROC). Essa curva ilustra a relação entre a taxa de verdadeiros positivos (sensibilidade) e a taxa de falsos positivos. Para resumir essa análise em um único valor, utiliza-se a área sob a curva (AUC), cujo valor varia de 0 a 1, sendo que valores próximos de 1 indicam maior capacidade discriminativa do modelo \cite{Faceli2011}.

Entretanto, conforme discutido por \cite{Davis2006}, a Curva ROC pode apresentar uma visão excessivamente otimista do desempenho em cenários com forte desbalanceamento de classes, como é comum em aplicações médicas. A Figura~\ref{fig:roc_pr_comparison} ilustra esse fenômeno: embora dois algoritmos apresentem desempenhos semelhantes no espaço ROC, diferenças significativas tornam-se evidentes quando analisados no espaço de Precisão-Revocação.

\begin{figure}[htpb]
\centering
\caption{Comparação entre Curvas ROC e Precisão-Revocação em dados desbalanceados}
\includegraphics[width=0.9\textwidth]{roc_pr.png}
\label{fig:roc_pr_comparison}
\fonte{Adaptado de \cite{Davis2006}}
\end{figure}

Essa discrepância ocorre porque a taxa de falsos positivos utilizada na Curva ROC é normalizada pelo número total de exemplos negativos, que tende a ser muito elevado em bases desbalanceadas, reduzindo o impacto visual dos erros. Em contraste, a Curva de Precisão-Revocação (PR) considera diretamente a proporção entre verdadeiros e falsos positivos, tornando-se mais sensível a variações no desempenho do modelo na classe minoritária.

Dessa forma, em cenários oncológicos, onde o interesse principal está na correta identificação de casos positivos (tumores malignos), a Curva PR e sua respectiva área sob a curva (PR-AUC) fornecem uma avaliação mais fiel do desempenho do classificador. Uma alta PR-AUC indica que o modelo consegue manter elevada precisão mesmo com altos níveis de sensibilidade, sendo, portanto, uma métrica mais adequada para validação em problemas com classes desbalanceadas.

\subsection{Principais Algoritmos Classificadores}

Com a base de dados metodologicamente preparada e as métricas de avaliação estabelecidas, o desenvolvimento do sistema avança para o seu núcleo computacional: a escolha do algoritmo de classificação. O campo do aprendizado supervisionado dispõe de um vasto leque de modelos matemáticos, variando desde equações probabilísticas mais tradicionais até arquiteturas estruturalmente complexas. Embora todos esses algoritmos compartilhem o mesmo propósito central de traçar uma fronteira capaz de separar as células benignas das malignas, cada um deles interpreta os dados e resolve esse problema geométrico utilizando lógicas completamente distintas.

O modelo de base mais clássico na modelagem estatística em saúde é a Regressão Logística. Apesar do nome referenciar problemas contínuos, trata-se de um classificador probabilístico. Ele atua calculando uma combinação linear das características da célula, como o seu raio e a sua textura, e aplica esse resultado numérico em uma função matemática chamada sigmoide. Essa função comprime qualquer valor para um intervalo estrito entre zero e um, fornecendo diretamente a probabilidade de o tumor ser maligno. Se essa probabilidade exceder um limiar de corte, o sistema consolida o diagnóstico positivo. A sua adoção contínua na prática clínica decorre da transparência da equação, permitindo que o médico entenda o peso exato de cada variável na decisão final \cite{Bishop2006}.

Quando a dispersão entre as características de células benignas e malignas apresenta uma distribuição complexa que não pode ser separada por uma linha reta simples, a computação recorre às Máquinas de Vetores de Suporte, frequentemente referenciadas pela sigla SVM. O objetivo primário deste algoritmo é encontrar uma fronteira de decisão geométrica, denominada hiperplano, que consiga separar as duas classes de tumores garantindo a maior margem de segurança possível entre elas. Os vetores de suporte, que dão nome ao modelo, são justamente os dados dos pacientes que estão posicionados de forma mais arriscada, ou seja, mais próximos a essa fronteira crítica. Eles atuam como os pilares estruturais que sustentam e definem o formato da decisão final do algoritmo \cite{Bishop2006}.

Entretanto, nos problemas oncológicos reais, é muito comum que os dados das células estejam sobrepostos, configurando um cenário não linear. Para solucionar essa limitação e viabilizar a classificação, o SVM utiliza uma operação matemática sofisticada conhecida como truque de \textit{kernel}. Conforme detalhado por \citeonline{AIOnlineCourse2023}, o \textit{kernel} atua como uma função de mapeamento que eleva os dados originais de um espaço de baixa dimensão para um espaço de dimensão superior. Nessa nova perspectiva matemática ampliada, as células malignas e benignas que antes pareciam misturadas passam a se afastar, tornando-se linearmente separáveis por um hiperplano. Existem diversas funções capazes de realizar essa transformação geométrica, sendo o \textit{kernel} Polinomial e o \textit{kernel} Gaussiano RBF os mais utilizados na prática para garantir que o algoritmo se adapte perfeitamente à complexidade biológica do câncer.

Outra abordagem geométrica muito intuitiva para a classificação é o algoritmo KNN(\textit{K-Nearest Neighbors}). Diferente dos modelos que processam os dados para construir uma equação fixa durante o treinamento, o KNN adota um método de aprendizado baseado em instâncias. Na prática, ele memoriza todo o conjunto de dados históricos de tumores. Quando os dados de um novo paciente são inseridos no sistema, o algoritmo calcula a distância matemática espacial, como a distância Euclidiana, entre as características dessa nova célula e todas as células antigas já armazenadas. O sistema então identifica os "K"\ vizinhos mais próximos, ou seja, os exames históricos que possuem as medidas anatômicas mais idênticas ao caso atual. O diagnóstico preditivo é determinado por uma votação simples nessa vizinhança. Se a maioria dos tumores vizinhos for maligna, o novo caso será classificado como maligno. A sua aplicação em sistemas médicos destaca-se pela facilidade de interpretação espacial por parte do profissional de saúde, embora o esforço computacional para calcular essas distâncias aumente consideravelmente em bases de dados muito volumosas \cite{Faceli2011}.

Explorando uma abordagem fundamentada em regras lógicas sucessivas, surgem os algoritmos baseados em Árvores de Decisão. Uma árvore individual funciona como um fluxograma hierárquico, testando condições passo a passo para classificar a célula oncológica. A matemática que define qual pergunta deve ser feita primeiro baseia-se frequentemente no critério de entropia, um conceito oriundo da teoria da informação que mede o grau de desordem ou impureza de um grupo de dados. Em cada etapa do fluxograma, o algoritmo calcula qual característica celular, como o raio ou a simetria, é capaz de dividir os pacientes de modo a reduzir ao máximo essa entropia. O objetivo é criar ramificações onde os tumores fiquem cada vez mais organizados e separados em suas classes puras. Contudo, quando crescem sem limites estritos, árvores de decisão isoladas tornam-se extremamente suscetíveis ao \textit{overfitting} \cite{Faceli2011}.

Para sanar essa vulnerabilidade da árvore única, a estatística computacional recorre à técnica \textit{Random Forest}, ou Floresta Aleatória. Este modelo opera sob o paradigma de comitês de máquinas, criando grandes quantidades de árvores de decisão simultâneas de forma paralela. Para garantir que as árvores não tomem decisões idênticas, o algoritmo força uma diversidade matemática ao treinar cada uma delas com um subconjunto aleatório de pacientes e características. Ao submeter os dados de um novo exame ao modelo, o diagnóstico final não é imposto por uma única regra, mas determinado por uma votação majoritária entre todas as árvores que compõem a floresta. Essa dinâmica de sabedoria coletiva dilui os erros de árvores individuais e gera um sistema preditivo altamente estável e protegido contra anomalias \cite{Faceli2011}.

Ainda dentro do paradigma de comitês de máquinas, destaca-se o \textit{Extreme Gradient Boosting}, amplamente conhecido no mercado como \textit{XGBoost}. Ao contrário da floresta aleatória que constrói todas as suas árvores de forma simultânea e independente, este algoritmo adota uma arquitetura sequencial focada no aprendizado com os próprios erros. A primeira árvore tenta classificar os tumores e inevitavelmente comete equívocos. A segunda árvore é construída com o propósito exclusivo de corrigir os erros da primeira, a terceira foca nos erros da segunda, e assim sucessivamente. Essa otimização contínua baseada na descida do gradiente cria um modelo preditivo de altíssima performance \cite{Faceli2011}.

No ápice da complexidade arquitetônica encontram-se as Redes Neurais Artificiais, algoritmos inspirados no funcionamento biológico do cérebro humano. Elas são compostas por múltiplas camadas de neurônios matemáticos interconectados, capazes de extrair relacionamentos profundos ocultos nos dados oncológicos \cite{Bishop2006}. Entretanto, essa evolução tecnológica impõe um grave dilema clínico. À medida que a ciência avança da transparente Regressão Logística para as poderosas redes neurais e algoritmos sequenciais como o \textit{XGBoost}, o modelo ganha acurácia estatística, mas perde drasticamente a sua interpretabilidade, transformando-se em uma caixa preta inescrutável. Na prática médica, entregar uma predição isolada não é eticamente aceitável, pois o oncologista necessita compreender as justificativas celulares do sistema. É exatamente para romper essa barreira analítica e resgatar a transparência das decisões que a Inteligência Artificial Explicável se torna indispensável.

\section{O Dilema da Caixa Preta e a Inteligência Artificial Explicável}

A evolução dos sistemas computacionais na oncologia transitou de ferramentas visuais puramente assistenciais para algoritmos preditivos autônomos de alta complexidade. Modelos baseados em comitês de árvores e redes neurais profundas alcançam taxas de acerto notáveis na classificação de tumores. Contudo, essa elevada acurácia estatística impôs um custo severo à prática clínica, consolidando o que a literatura científica denomina como o dilema da caixa preta. Em um cenário de alto risco como o diagnóstico de câncer de mama, a previsão matemática isolada é insuficiente. O médico patologista não pode basear uma decisão terapêutica drástica, como uma intervenção cirúrgica, em um sistema que emite um alerta de malignidade sem fornecer o raciocínio biológico e morfológico que motivou tal conclusão \cite{Holzinger2019}.

Para preencher essa lacuna entre a precisão matemática e a confiança médica, consolidou-se o campo da Inteligência Artificial Explicável, amplamente referenciado pela sigla em inglês XAI. O propósito primordial da XAI não é focar na criação de novos classificadores, mas sim no desenvolvimento de métodos e ferramentas que tornem o processo de decisão das máquinas transparente e compreensível para os seres humanos. Em aplicações médicas, a explicabilidade atua como uma camada de auditoria e segurança. Ela permite que o especialista de saúde verifique se o algoritmo está baseando sua decisão em características celulares clinicamente válidas, como a concavidade e a textura do núcleo, em vez de aprender padrões espúrios ou artefatos irrelevantes presentes na base de dados \cite{Adadi2018}.

Do ponto de vista arquitetônico, \citeonline{Rudin2019} destaca que a explicabilidade se divide em duas abordagens fundamentais. A primeira consiste na utilização de modelos intrinsecamente interpretáveis, como a Regressão Logística ou uma única Árvore de Decisão, cujas equações matemáticas são naturalmente transparentes, embora frequentemente apresentem limitações de desempenho em dados biológicos complexos. A segunda abordagem engloba as técnicas de explicabilidade \textit{post-hoc}. Estas técnicas são acopladas externamente a modelos de caixa preta após a fase de treinamento, extraindo interpretações sem alterar a estrutura interna do algoritmo original. No escopo deste estudo, opta-se pela segunda via. Ao aplicar métodos \textit{post-hoc} sobre classificadores robustos, busca-se entregar ao oncologista uma ferramenta que alie o máximo desempenho preditivo da computação moderna à transparência analítica exigida pela ética médica.

\section{SHapley Additive exPlanations (SHAP)}

A técnica \textit{SHapley Additive exPlanations} (SHAP) representa o atual estado da arte em explicabilidade algorítmica. A sua base teórica não nasceu na ciência da computação, mas sim na teoria dos jogos cooperativos, formulada pelo matemático Lloyd Shapley. Na teoria original, o valor de Shapley calcula a contribuição exata e justa de cada jogador para o prêmio final obtido por uma equipe. Quando esse conceito é transportado para o aprendizado de máquina, o modelo preditivo assume o papel do jogo, as características celulares atuam como os jogadores e o prêmio a ser dividido é a predição final do algoritmo. De acordo com a documentação oficial da ferramenta, o SHAP atua como uma abordagem unificada, consolidando diversas metodologias anteriores sob a classe de métodos de atribuição de características aditivas para revelar os bastidores matemáticos da caixa preta \cite{Lundberg2017}.

Um dos grandes trunfos arquitetônicos da biblioteca SHAP, conforme detalhado em seu manual técnico, é a disponibilidade de otimizadores algorítmicos específicos denominados \textit{Explainers}. Como o cálculo de todas as combinações possíveis de variáveis possui um custo computacional altíssimo, a ferramenta oferece soluções otimizadas para diferentes modelos. Para algoritmos baseados em árvores, como o \textit{Random Forest} e o \textit{XGBoost}, a biblioteca disponibiliza o \textit{TreeExplainer}, um algoritmo capaz de calcular os valores exatos de Shapley de forma extremamente rápida aproveitando a estrutura de roteamento interna das árvores. Para arquiteturas matemáticas distintas ou caixas pretas absolutas, o sistema recorre a algoritmos agnósticos, como o \textit{KernelExplainer}, garantindo que qualquer modelo oncológico possa ser dissecado independentemente de sua complexidade teórica \cite{SHAPDocs2026}.

O resultado desse robusto processamento matemático é traduzido em visualizações gráficas projetadas especificamente para a prática clínica e auditoria de dados. No escopo da interpretabilidade local, voltada para o diagnóstico de um único paciente, o SHAP fornece gráficos de cascata (\textit{waterfall plots}). Essas representações ilustram milimetricamente como cada característica, como a textura ou a área do núcleo celular, empurrou a probabilidade do diagnóstico para cima ou para baixo a partir do valor base da população. No escopo global, o algoritmo compila todas as predições em gráficos de enxame (\textit{beeswarm plots} ou \textit{summary plots}). Esse panorama permite que o médico observe não apenas quais variáveis são globalmente mais importantes para o classificador, mas também a direção biomédica desse impacto, revelando se os valores altos de uma determinada característica estão sistematicamente associados a tumores malignos.

Em comparação com outras ferramentas de explicabilidade disponíveis na literatura, como o modelo LIME, o método SHAP destaca-se por sua extrema robustez teórica. Enquanto outras técnicas tentam adivinhar o comportamento do algoritmo criando aproximações em torno dos dados locais, o SHAP garante propriedades matemáticas vitais de consistência e precisão local. Na prática, isso significa que a soma de todos os valores e pesos atribuídos a cada característica celular pelo SHAP resultará exatamente na predição percentual final emitida pelo sistema principal. Essa garantia de exatidão torna o método estatisticamente seguro e incontestável para auditar diagnósticos críticos na área da saúde.

\subsection{Fundamentação Matemática: A Equação de Shapley}

Para compreender como a biblioteca computacional distribui os pesos das características oncológicas, é estritamente necessário analisar a formulação matemática original proposta para o Valor de Shapley. O cálculo baseia-se na análise de todas as permutações e subconjuntos possíveis de jogadores em um cenário cooperativo. Seja $N$ o conjunto total de características celulares disponíveis em um exame e $S$ um subconjunto específico dessas características, o valor de Shapley para uma característica isolada $i$, denotado por $\phi_i$, é definido pela Equação \ref{eq:shapley} \cite{Lundberg2017}:

\begin{equation}
    \label{eq:shapley}
    \phi_i = \sum_{S \subseteq N \setminus \{i\}} \frac{|S|! (|N| - |S| - 1)!}{|N|!} [v(S \cup \{i\}) - v(S)]
\end{equation}

A interpretação desta equação revela a essência do método explicativo. O termo $v(S \cup \{i\}) - v(S)$ representa a contribuição marginal da característica $i$. Na prática médica, ele calcula a diferença entre a predição de risco do modelo quando a característica $i$ (como a simetria celular) está presente no exame e a predição quando ela é omitida daquela combinação específica. O algoritmo calcula essa diferença não apenas uma vez, mas para todos os subconjuntos $S$ possíveis formados pelas outras variáveis do modelo.

A fração que antecede a contribuição marginal atua como um sistema de pesos baseado em probabilidade combinatória. Os termos fatoriais, representados pelo símbolo de exclamação, calculam o peso exato daquela coalizão específica se formar, garantindo que o valor final seja distribuído de maneira perfeitamente justa. Por fim, o somatório agrega todas essas contribuições marginais ponderadas. O resultado é um único valor $\phi_i$ que quantifica, com exatidão matemática, o impacto líquido que aquela característica anatômica exerceu para elevar ou reduzir a probabilidade de o tumor ser classificado como maligno.

\section{Redução de Dimensionalidade e UMAP}

A análise preditiva de tumores lida com um volume massivo de informações simultâneas. No caso do câncer de mama, o algoritmo avalia dezenas de características anatômicas em um único exame, como o raio, a textura e a simetria de cada núcleo celular \cite{Street1993}. Embora os modelos matemáticos consigam processar e cruzar esses dados em um espaço de trinta dimensões com extrema facilidade, essa estrutura de alta dimensionalidade não possui uma representação física direta que possa ser desenhada ou interpretada visualmente. Essa assimetria estrutural cria um obstáculo analítico significativo na prática médica, pois inviabiliza a observação direta do comportamento de toda a base de pacientes utilizando os eixos convencionais de um gráfico bi ou tridimensional.

Para solucionar essa limitação biomédica, a ciência de dados recorre a técnicas de redução de dimensionalidade. O propósito primário dessas ferramentas não é simplesmente descartar variáveis anatômicas aleatoriamente, mas sim aplicar transformações matemáticas capazes de comprimir um grande volume de dimensões em apenas dois eixos principais. O desafio central desse processo de compressão é preservar ao máximo a essência da informação original. Uma técnica metodologicamente eficiente deve garantir que a estrutura dos dados apresentada no gráfico bidimensional seja um reflexo direto e fiel das relações biológicas que existiam no espaço multidimensional complexo.

Entre as metodologias disponíveis na literatura estatística contemporânea, o algoritmo UMAP, que consiste na Aproximação e Projeção de Variedade Uniforme, consolidou-se como uma ferramenta de altíssimo desempenho para a redução de dimensionalidade. O seu funcionamento interno diverge de abordagens matemáticas puramente lineares ao fundamentar-se em princípios de topologia algébrica. Na primeira etapa de seu processamento computacional, o algoritmo constrói uma representação teórica do conjunto de dados em seu espaço original de alta dimensão. Para isso, o modelo calcula as distâncias entre todos os pontos do sistema e mapeia rigorosamente os vizinhos mais próximos de cada elemento, formando um complexo grafo matemático que modela a estrutura exata de conectividade da informação original \cite{McInnes2018}.

Uma vez que essa rede de vizinhanças é estabelecida, a segunda fase do algoritmo consiste em projetar essa estrutura em um espaço de dimensão reduzida, como um plano bidimensional. Para realizar essa transição geométrica sem descaracterizar os dados, o UMAP inicializa um posicionamento espacial e aplica um rigoroso processo de otimização. O sistema utiliza vetores de atração matemática para aproximar as coordenadas dos pontos que foram mapeados como vizinhos fortes no espaço original, ao mesmo tempo em que aplica métricas de repulsão para afastar os pontos que não possuíam conexões de similaridade. O algoritmo ajusta iterativamente essas posições até que o gráfico plano atinja a aproximação topológica mais fiel possível da complexa rede inicial.

O grande diferencial técnico desta arquitetura algorítmica reside na sua capacidade matemática de preservar simultaneamente a estrutura local e a estrutura global do conjunto de dados originais. Ao contrário de técnicas anteriores que frequentemente sacrificam a disposição geral em prol do ajuste de vizinhanças imediatas, o UMAP assegura que as proporções de distância sejam mantidas na projeção visual. Como resultado geométrico direto desse processamento topológico, elementos de dados com propriedades numéricas semelhantes convergem naturalmente para regiões isoladas do gráfico, formando densos agrupamentos visuais conhecidos na estatística descritiva como \textit{clusters}. Essa capacidade inerente de revelar a separabilidade natural de bases complexas consolida o método como um recurso indispensável na análise exploratória avançada.






% % -------------------------------------------------------------------------------------------------
% % -------------------------------------------------------------------------------------------------
% % 3. TRABALHOS RELACIONADOS
% % -------------------------------------------------------------------------------------------------
% % -------------------------------------------------------------------------------------------------

\chapter{Revisão da Literatura}

Este capítulo apresenta uma análise detalhada da literatura científica que fundamenta o estado da arte no uso de aprendizado de máquina para o diagnóstico oncológico. A seleção abrange desde os estudos seminais que deram origem às bases de dados modernas até pesquisas recentes focadas em inteligência artificial explicável, estabelecendo o contexto necessário para a identificação da lacuna de pesquisa preenchida por este projeto.

\section{Extração de Características Nucleares para o Diagnóstico de Tumores Mamários (Street et al., 1993)}

O trabalho conduzido por \citeonline{Street1993} é considerado a pedra angular para o desenvolvimento de sistemas de diagnóstico assistido por computador aplicados ao câncer de mama. O objetivo primordial da pesquisa foi desenvolver um sistema que pudesse processar imagens digitais de aspirados por agulha fina para fornecer um diagnóstico maligno ou benigno de forma automatizada e precisa.

A metodologia baseou-se na criação de um sistema de análise de imagens que permitia ao usuário delimitar grosseiramente as bordas dos núcleos celulares. A partir dessa interação inicial, o software calculava automaticamente dez características geométricas distintas para cada núcleo, incluindo o raio, o perímetro, a textura e a simetria. Esse processo gerou a base de dados conhecida como \textit{Wisconsin Diagnostic Breast Cancer} (WDBC), que é utilizada como padrão ouro na literatura acadêmica até os dias atuais.

Os resultados obtidos na época demonstraram uma acurácia de classificação de noventa e sete por cento utilizando um classificador baseado em programação linear. O trabalho provou que as medidas nucleares, quando extraídas de forma sistemática, continham informações estatísticas suficientes para diferenciar tumores sem a necessidade de intervenções invasivas imediatas.

Para este projeto de graduação, este estudo é fundamental pois define a origem e o significado biológico de cada uma das variáveis que serão processadas pelos modelos de SHAP e UMAP. Compreender a gênese dos dados permite uma interpretação mais rica das explicações fornecidas pela inteligência artificial explicável.

\section{Aplicações de Machine Learning na Predição e Prognóstico do Câncer (Cruz e Wishart, 2006)}

A revisão sistemática elaborada por \citeonline{Cruz2006} ofereceu um panorama crítico sobre como as técnicas de aprendizado de máquina estavam migrando de laboratórios de pesquisa para aplicações práticas no prognóstico oncológico. O foco do estudo foi avaliar a viabilidade de substituir métodos estatísticos convencionais por modelos preditivos mais robustos.

Os autores analisaram diversas arquiteturas, como as Máquinas de Vetores de Suporte (SVM) e Redes Neurais Artificiais, destacando que a complexidade biológica do câncer exigia ferramentas capazes de lidar com relacionamentos não lineares. A pesquisa ressaltou que, embora o câncer de mama fosse o mais estudado, a metodologia poderia ser estendida para outros tipos de neoplasias, como as de próstata e pulmão.

Uma das conclusões mais relevantes deste trabalho foi a identificação de que a acurácia isolada não era o único fator de sucesso. Os pesquisadores apontaram que a qualidade dos dados de entrada e a seleção cuidadosa de variáveis eram determinantes para a utilidade clínica do modelo.

Este estudo corrobora a necessidade de superar as limitações da estatística básica na oncologia. No entanto, a busca por capturar as nuances da malignidade celular não deve resultar em modelos de complexidade interpretativa inacessível. A abordagem deste trabalho enfatiza que, no contexto médico, modelos preditivos que priorizam a simplicidade e a explicabilidade trazem vantagens substanciais, permitindo que a equipe de saúde compreenda e confie nos critérios lógicos que fundamentam o diagnóstico.

\section{Segmentação e Contagem de Células em Imagens (Silva Júnior, 2018)}

Desenvolvido no contexto acadêmico de Curitiba, o trabalho de \citeonline{SilvaJunior2018} focou no desenvolvimento de um \textit{pipeline} detalhado para o processamento digital de imagens microscópicas de câncer de mama. O objetivo era automatizar a tarefa exaustiva de contagem e segmentação de núcleos celulares realizada por patologistas.

A abordagem metodológica utilizou técnicas de processamento de imagens para isolar componentes celulares e reduzir o ruído visual das lâminas. O autor detalhou cada etapa do fluxo de trabalho, desde a captura da imagem até a extração final de dados quantitativos, enfatizando os desafios técnicos de lidar com células sobrepostas ou imagens de baixa qualidade.

Embora o foco principal fosse a segmentação e não a classificação preditiva, o trabalho demonstrou a importância de transformar evidências visuais em dados tabulares estruturados. O autor concluiu que a automação dessas etapas iniciais reduz significativamente a fadiga do profissional de saúde e aumenta a reprodutibilidade dos diagnósticos.

A relevância deste estudo para o presente projeto reside na compreensão da infraestrutura necessária para gerar os dados de entrada. Enquanto este TCC foca na inteligência preditiva e explicável sobre dados já estruturados, o trabalho de Silva Júnior fornece a base sobre como esses dados são extraídos do mundo físico.

\section{Diagnóstico de Câncer de Pulmão e Cólon via Deep Learning (Masud et al., 2021)}

A pesquisa de \citeonline{Masud2021} explorou o uso de Redes Neurais Profundas para o diagnóstico de câncer em tecidos pulmonares e colorretais. O diferencial do trabalho foi o uso de uma estrutura de classificação baseada em imagens de histopatologia para identificar padrões malignos de forma automática.

O modelo desenvolvido utilizou técnicas de aprendizado profundo para extrair características de alto nível que muitas vezes são imperceptíveis ao olho humano. Os autores conseguiram alcançar métricas de desempenho superiores a noventa e seis por cento de acurácia, demonstrando o poder das arquiteturas modernas de inteligência artificial na oncologia.

Contudo, os pesquisadores notaram que modelos de rede neural profunda, apesar de precisos, operam como sistemas de difícil interpretação. O estudo concluiu que a integração de sistemas computacionais na medicina requer um equilíbrio entre a capacidade de processamento e a facilidade de uso pelo corpo clínico.

Este trabalho é citado para contrastar com a abordagem tabular deste TCC. Ele exemplifica o estado da arte em processamento de imagens, servindo de base para discutir por que a explicabilidade se torna ainda mais necessária quando lidamos com modelos de alta complexidade como os de \textit{Deep Learning}.

\section{Predição de Câncer de Pulmão com Seleção de Correlação (Abdullah et al., 2021)}

Focando na eficiência computacional, \citeonline{Abdullah2021} investigaram como a seleção inteligente de variáveis pode melhorar a precisão do diagnóstico de câncer de pulmão. O objetivo era identificar quais características clínicas e genéticas eram mais correlacionadas com a presença da doença.

A metodologia utilizou a ferramenta WEKA para testar diversos classificadores, com destaque para as Máquinas de Vetores de Suporte (SVM). Os autores aplicaram um método de seleção baseado em correlação para remover variáveis irrelevantes ou redundantes, buscando um modelo mais enxuto e rápido.

Os resultados mostraram que modelos com menos variáveis, porém mais significativas, apresentavam desempenho equivalente ou superior a modelos que utilizavam todos os dados disponíveis. Isso demonstrou que o excesso de informação pode, por vezes, prejudicar o aprendizado da máquina.

Este estudo corrobora a utilidade do UMAP no presente projeto. Ao discutir a redução de dimensionalidade, o trabalho de Abdullah et al. fornece o embasamento sobre como a simplificação de dados complexos pode levar a diagnósticos mais precisos e interpretáveis.

\section{Abordagens de Machine Learning na Predição do Câncer de Mama (Rabiei et al., 2022)}

O estudo de \citeonline{Rabiei2022} avaliou diferentes algoritmos de aprendizado de máquina aplicados a dados hospitalares reais de câncer de mama. O foco principal foi lidar com o problema de bases de dados desbalanceadas, onde o número de casos benignos costuma ser muito superior aos malignos.

Os pesquisadores testaram modelos como Regressão Logística, \textit{Random Forest} e Redes Neurais. Para corrigir o desbalanceamento, utilizaram técnicas de sobreamostragem, permitindo que os algoritmos aprendessem com a mesma eficácia sobre ambas as classes de tumores.

As conclusões apontaram que o algoritmo \textit{Random Forest} obteve o melhor equilíbrio entre sensibilidade e especificidade. O trabalho enfatizou que o ajuste de hiperparâmetros e o pré-processamento rigoroso são as chaves para a obtenção de modelos confiáveis em ambientes médicos.

Para este TCC, a pesquisa de Rabiei et al. serve como uma referência metodológica para a fase de treinamento dos modelos. Ela valida a inclusão de comitês de árvores no conjunto de classificadores a serem testados e explicados via SHAP.

\section{Predição Baseada em Aprendizado de Máquina (Chen et al., 2023)}

A pesquisa de \citeonline{Chen2023} revisitou o desafio da classificação do câncer de mama focando na robustez das métricas de avaliação. O objetivo era determinar quais algoritmos eram mais seguros para a adoção clínica, priorizando a redução de falsos negativos.

O estudo utilizou a base de dados de Wisconsin para treinar modelos e compará-los através de curvas ROC e métricas de precisão. Os autores destacaram que a inteligência artificial não deve ser vista como um substituto para o médico, mas como uma ferramenta de segunda opinião capaz de filtrar casos suspeitos com alta velocidade.

Os resultados confirmaram a eficácia dos modelos de \textit{Boosting} para este tipo de dado tabular. A pesquisa concluiu que a integração desses modelos em sistemas de apoio à decisão poderia acelerar significativamente o início do tratamento de pacientes oncológicos.

Este trabalho é relevante por ser uma das publicações mais recentes sobre a mesma base de dados utilizada neste projeto. Ele estabelece o patamar de acurácia que este TCC deve buscar atingir ou superar antes de aplicar as camadas de explicabilidade.

\section{Estudo Comparativo com Seleção LASSO e SHAP (Shaon et al., 2024)}

Considerado o trabalho mais próximo da proposta deste TCC, \citeonline{Shaon2024} realizaram uma comparação de modelos preditivos utilizando técnicas modernas de seleção de atributos e explicabilidade. O objetivo era não apenas prever o câncer, mas entender quais variáveis eram determinantes para o modelo.

A metodologia integrou o uso de LASSO para a regularização dos dados e o SHAP para a interpretabilidade local e global. Os autores alcançaram uma acurácia de noventa e nove por cento, demonstrando que é possível unir alto desempenho com transparência analítica.

O estudo concluiu que o uso de valores de Shapley oferece uma justificativa robusta para cada diagnóstico, permitindo que médicos confiem nos resultados apresentados. No entanto, o trabalho limitou-se à análise técnica dos dados, sem propor uma ferramenta de interface para o usuário final.

Este artigo representa o "estado da arte"\ atual e serve como o principal ponto de comparação. O diferencial deste TCC em relação ao trabalho de Shaon et al. é a inclusão do UMAP para visualização topológica e o desenvolvimento de um software completo com interface gráfica, indo além da análise estatística pura.

\section{Síntese e Lacuna de Pesquisa}

A análise dos trabalhos relacionados permite concluir que, embora existam pesquisas robustas focadas na acurácia preditiva e estudos recentes explorando a técnica SHAP, ainda há uma carência de soluções que integrem esses elementos em uma ferramenta de software coesa e utilizável pelo profissional de saúde.

Muitos dos modelos propostos na literatura permanecem restritos a ambientes de programação ou relatórios estáticos. O Quadro \ref{tab:trabalhos_relacionados} sintetiza a comparação entre as obras citadas e a proposta deste projeto, evidenciando que a união de classificação de alta performance, explicabilidade (SHAP), redução de dimensionalidade visual (UMAP) e interface de processamento em lote constitui o diferencial competitivo e a inovação acadêmica deste trabalho.

\begin{table}[ht]
    \centering
    \caption{Comparação entre Trabalhos Relacionados e o Presente Projeto}
    \label{tab:trabalhos_relacionados}
    \resizebox{\textwidth}{!}{
    \begin{tabular}{lcccccc}
        \toprule
        Trabalho & ML & Base Pública & Tabular & Mama & Interface & XAI \\
        \midrule
        Street et al. (1993) & Sim & Não & Não & Sim & Sim & Não \\
        Cruz \& Wishart (2006) & Sim & Sim & Sim & Não & Não & Não \\
        Silva Júnior (2018) & Não & Sim & Não & Sim & Não & Não \\
        Masud et al. (2021) & Sim & Sim & Não & Não & Não & Não \\
        Abdullah (2021) & Sim & Sim & Sim & Não & Não & Não \\
        Rabiei et al. (2022) & Sim & Não & Sim & Sim & Não & Não \\
        Chen et al. (2023) & Sim & Sim & Sim & Sim & Não & Não \\
        Shaon et al. (2024) & Sim & Sim & Sim & Sim & Não & Sim \\
        \midrule
        Este Projeto & Sim & Sim & Sim & Sim & Sim & Sim \\
        \bottomrule
    \end{tabular}
    }
\end{table}






% % -------------------------------------------------------------------------------------------------
% % -------------------------------------------------------------------------------------------------
% % 2. ORGANIZAÇÃO
% % -------------------------------------------------------------------------------------------------
% % -------------------------------------------------------------------------------------------------
% \chapter{Organização do Trabalho}
% \label{cap:organizacao}

% Este capítulo apresenta as principais seções de um trabalho acadêmico, discute a hierarquia de seções  e fornece exemplos de como implementar a sua organização. A organização de um trabalho acadêmico é um ponto fundamental para garantir a clareza e a estrutura lógica do texto~\cite{severino}.

% Normalmente, um trabalho acadêmico possui a seguinte divisão de capítulos:

% \begin{alineas}
%     \item Introdução: estabelece o contexto do trabalho e apresenta o problema de pesquisa;
%     \item Fundamentação teórica: explora os conceitos e teorias que embasam a pesquisa, proporcionando uma base teórica para a compreensão do tema;
%     \item Revisão da Literatura: analisa e discute estudos anteriores relacionados ao tema, identificando lacunas que a pesquisa busca preencher;
%     \item Metodologia: descreve os métodos e técnicas utilizados para conduzir a pesquisa, permitindo a reprodução do estudo por outros pesquisadores;
%     \item Resultados Experimentais: apresenta e analisa os dados obtidos durante a pesquisa, evidenciando os achados e sua relevância para o tema investigado;
%     \item Conclusão: resume os principais resultados, discute suas implicações, e sugere direções para futuras pesquisas ou aplicações práticas.
% \end{alineas}

% A organização adequada de um trabalho acadêmico não apenas melhora a legibilidade, mas também facilita a navegação pelo documento. O uso correto dos comandos de seção, subseção e subsubseção ajuda a estruturar o conteúdo de forma lógica e coesa, orientando o leitor através dos diferentes tópicos abordados. Além disso, a numeração automática das seções e subseções torna o trabalho mais profissional e alinhado às normas acadêmicas~\cite{gil}.



% % -------------------------------------------------------------------------------------------------
% % -------------------------------------------------------------------------------------------------
% % -------------------------------------------------------------------------------------------------


% % -------------------------------------------------------------------------------------------------
% % -------------------------------------------------------------------------------------------------
% % 3. CITAÇÃO
% % -------------------------------------------------------------------------------------------------
% % -------------------------------------------------------------------------------------------------
% \chapter{Citação}
% \label{cap:citacao}

% Este capítulo apresenta os principais tipos de citação empregados em trabalhos acadêmicos, detalhando suas características e a importância de sua correta aplicação para fundamentar teoricamente a pesquisa. Citações diretas e indiretas são essenciais para a construção de uma argumentação sólida, e o uso de ferramentas como o BibTeX facilita a organização e padronização das referências bibliográficas. Ao final, são discutidas as melhores práticas para citação de livros, artigos, capítulos de livros, teses, dissertações e fontes online, com vistas à preservação da integridade acadêmica.

% \section{Tipos de Citações}

% As citações são classificadas em dois tipos principais: (1) diretas e (2) indiretas, sendo ambas complementares no desenvolvimento da argumentação e fundamentação teórica de um trabalho.

% \subsection{Citação Direta}

% A citação direta consiste na reprodução literal das palavras de um autor. Segundo Gil~\cite{gil}, ela é utilizada quando se deseja destacar uma formulação precisa ou quando o uso exato das palavras é fundamental para a compreensão do argumento. No entanto, seu uso deve ser comedido e sempre referenciado adequadamente.

% As citações diretas curtas podem ser inseridas com aspas e acompanhadas da devida referência, como demonstrado: ``O conhecimento científico é sistemático e controlado''~\cite{gil}. Quando a citação excede três linhas, deve ser formatada em bloco, sem aspas, como no exemplo a seguir:

% \begin{citacao}
%     O conhecimento científico é um conhecimento crítico, que tenta resolver os problemas que surgem em sua busca pela compreensão dos fenômenos, submetendo suas hipóteses e teorias a rigorosos testes de falseabilidade~\cite{gil}.
% \end{citacao}

% Por se tratar de uma reprodução exata, a citação direta deve preservar a ortografia, a pontuação e o uso de maiúsculas do texto original. Qualquer alteração ou omissão no texto original deve ser indicada com colchetes. Mesmo os erros presentes no texto original não devem ser corrigidos~\cite{utp_normas2024}. Para isso, algumas convenções são seguidas, como:

% \begin{alineas}
%     \item {[...]}: indica a omissão de parte do texto original;
%     \item {[sic]}: sinaliza um erro ou incoerência no texto original, significando "assim mesmo";
%     \item {[grifo meu]}, [grifo nosso], [sem grifo no original]: utilizados quando o autor que cita adiciona ênfase (negrito, sublinhado, etc.) que não estava no texto original;
%     \item {[grifo do autor]}: indica que o grifo já fazia parte do texto original citado;
%     \item Alterações ou comentários feitos na citação direta devem ser incluídos entre colchetes: [ ].
% \end{alineas}

% Citações em idiomas estrangeiros devem ser traduzidas, acompanhadas da devida referência. Quando a tradução for feita pelo próprio autor do trabalho, deve-se incluir a nota ``tradução do autor(a)'' e a versão original pode ser inserida como nota de rodapé. No entanto, em textos de caráter literário, é comum manter o idioma original, especialmente quando a tradução poderia comprometer o estilo ou a análise literária~\cite{leite, utp_normas2024}.

% \subsection{Citação Indireta}

% A citação indireta ocorre quando as ideias de outro autor são parafraseadas ou resumidas. Essa prática permite ao autor integrar diferentes perspectivas em seu próprio raciocínio. Segundo \citeonline{creswell}, a citação indireta é uma ferramenta importante para articular múltiplas fontes e construir uma análise mais crítica. As citações indiretas podem ser feitas de duas maneiras: (1)~explícitas ou (2)~implícitas. 

% No caso da citação explícita, o nome do autor é mencionado no texto usando o comando \texttt{\textbackslash citeonline}, como no exemplo a seguir. De acordo com~\citeonline{creswell}, a pesquisa qualitativa depende de um processo interpretativo contínuo.

% Já na citação indireta implícita, as ideias de um ou mais autores são sintetizadas e as citações são incluídas ao final da frase ou parágrafo usando o comando \texttt{\textbackslash cite}. as referências são inseridas ao final do parágrafo ou frase, sem menção direta ao autor no texto, como a seguir. A pesquisa qualitativa envolve um processo de reflexão constante e interpretação dos dados~\cite{creswell}.

% \subsection{Citação de Fonte}

% Ao utilizar figuras ou tabelas de outras fontes, ou até mesmo do próprio autor, é fundamental dar o devido crédito à fonte original. Isso garante o reconhecimento adequado do trabalho original e evita problemas de plágio. A citação da fonte de figuras e tabelas pode ser feita de maneira semelhante à citação de textos, usando o comando \texttt{\textbackslash fonte}. A Figura~\ref{fig:fontes} ilustra como as fontes devem ser citadas~\cite{utp_normas2024}.

% \begin{figure}[htb]
%     \legenda[fig:fontes]{Exemplo de citação de fontes}
%     \fig{width=0.7\textwidth}{img/fontes}
%     \fonte{\citeauthoronline{utp_normas2024}, \citeyear{utp_normas2024}, p. 23}
% \end{figure}

% O Código~\ref{cod:fontes} exemplifica como que os diferentes tipos de fontes devem ser incluídos no código para adequação às normas técnicas.

% \begin{codigo}[htb]
% \legenda[cod:fontes]{Diferentes tipos de citação em fontes}
% \begin{lstlisting}
% % Citação de autor
% \fonte{\citeauthoronline{autor}, \citeyear{autor}, p. 99}

% % Citação de autor-entidade
% \fonte{\citeauthoronline{entidade}, \citeyear{entidade}, p. 105}

% % Citação de site
% \fonte{\citeauthoronline{site}, \citeyear{site}, Disponível em: \url{...}}

% % Citação de autoria própria 
% \fonteautor     % ou \fonteautora
% \fonteautores   % ou \fonteautoras
% \fonteautoria
% \end{lstlisting}
% \fonteautor
% \end{codigo}

% \section{Diferentes Tipos de Referências}

% Além de conhecer os tipos de citação, é importante dominar os diversos formatos de referências utilizados em diferentes tipos de publicações. O uso do BibTeX permite automatizar a formatação dessas referências, garantindo padronização. A seguir, são apresentados os formatos para as principais fontes acadêmicas.

% \subsection{Livros}

% Referências de livros são amplamente utilizadas na fundamentação teórica, pois oferecem uma visão geral e detalhada sobre os conceitos e teorias fundamentais de uma área de estudo. Elas são essenciais para embasar o desenvolvimento dos argumentos do autor. A referência de livro pode ser inserida no arquivo BibTeX conforme apresentado no Código~\ref{cod:bibtex-book}.

% \begin{codigo}[htb]
% \legenda[cod:bibtex-book]{Estrutura BibTeX para citação de um livro}
% \begin{lstlisting}
% @book{autor,
%     author    = {Sobrenome, Nome},
%     title     = {Título do Livro},
%     edition   = {Edição},
%     publisher = {Editora},
%     year      = {Ano},
%     pages     = {Número de Paginas}
%     address   = {Local de Publicacao}
% }
% \end{lstlisting}
% \fonteautor
% \end{codigo}

% \subsection{Artigos de Periódicos}

% Artigos científicos publicados em periódicos são uma fonte de informação atualizada e focada em tópicos específicos. Em um trabalho acadêmico, esses artigos são essenciais para manter a pesquisa em sintonia com as discussões mais recentes na área, e, por isso, sua correta citação é de extrema importância. A entrada BibTeX para artigos de periódicos é descrita no Código~\ref{cod:bibtex-article}.

% \begin{codigo}[htb]
% \legenda[cod:bibtex-article]{Estrutura BibTeX para citação de um artigo de periódico}
% \begin{lstlisting}
% @article{autor,
%     author  = {Sobrenome, Nome},
%     title   = {Título do Artigo},
%     journal = {Nome do Periódico},
%     volume  = {Volume},
%     number  = {Número},
%     pages   = {Páginas},
%     year    = {Ano}
% }
% \end{lstlisting}
% \fonteautor
% \end{codigo}

% \subsection{Capítulos de Livros}

% Muitas vezes, é necessário citar apenas um capítulo de um livro organizado por diversos autores. Isso ocorre quando o capítulo em questão é escrito por um autor diferente do editor do livro, ou quando se deseja dar destaque a uma seção específica da obra. Isto pode ser feito utilizando a estrutura apresentada no Código~\ref{cod:bibtex-chapter}.

% \begin{codigo}[htb]
% \legenda[cod:bibtex-chapter]{Estrutura BibTeX para citação de um capítulo de livro}
% \begin{lstlisting}
% @incollection{autor,
%     author    = {Sobrenome, Nome},
%     title     = {Título do Capítulo},
%     booktitle = {Título do Livro},
%     publisher = {Editora},
%     year      = {Ano},
%     address   = {Local de Publicação},
%     pages     = {Páginas do Capítulo}
% }
% \end{lstlisting}
% \fonteautor
% \end{codigo}

% \subsection{Tese de Doutorado}

% Teses de doutorado são referências relevantes em muitas áreas, especialmente porque contêm contribuições originais de pesquisa. Elas são uma fonte importante de conhecimento atualizado e inédito. A estrutura de citação para uma tese de doutorado em BibTeX pode ser vista no Código~\ref{cod:bibtex-phdthesis}.

% \begin{codigo}[htb]
% \legenda[cod:bibtex-phdthesis]{Estrutura BibTeX para citação de uma tese de doutorado}
% \begin{lstlisting}
% @phdthesis{autor,
%     author  = {Sobrenome, Nome},
%     title   = {Título da Tese},
%     school  = {Nome da Universidade},
%     year    = {Ano},
%     address = {Local da Universidade}
% }
% \end{lstlisting}
% \fonteautor
% \end{codigo}

% \subsection{Dissertação de Mestrado}

% As dissertações de mestrado, assim como as teses de doutorado, são fundamentais para a pesquisa científica, por apresentarem investigações originais ou revisões aprofundadas de temas específicos. A estrutura para citar uma dissertação de mestrado é semelhante à de teses, como mostrado no Código~\ref{cod:bibtex-mastersthesis}.

% \begin{codigo}[htb]
% \legenda[cod:bibtex-mastersthesis]{Estrutura BibTeX para citação de uma dissertação de mestrado}
% \begin{lstlisting}
% @mastersthesis{autor,
%     author  = {Sobrenome, Nome},
%     title   = {Título da Dissertaçãoo},
%     school  = {Nome da Universidade},
%     year    = {Ano},
%     address = {Local da Universidade},
%     type    = {Dissertação de Mestrado}
% }
% \end{lstlisting}
% \fonteautor
% \end{codigo}

% \subsection{Monografia de Graduação}

% Monografias de graduação são trabalhos acadêmicos exigidos para a conclusão de cursos de graduação e podem ser relevantes como base de estudos introdutórios ou revisões bibliográficas. Para citá-las, o formato é apresentado no Código~\ref{cod:bibtex-monography} e é similar à de dissertações.

% \begin{codigo}[htb]
% \legenda[cod:bibtex-monography]{Estrutura BibTeX para citação de uma monografia de graduação}
% \begin{lstlisting}
% @mastersthesis{autor,
%     author  = {Sobrenome, Nome},
%     title   = {Título da Monografia},
%     school  = {Nome da Universidade},
%     year    = {Ano},
%     address = {Local da Universidade},
%     type    = {Monografia de Graduação}
% }
% \end{lstlisting}
% \fonteautor
% \end{codigo}

% \subsection{Conferências}

% Artigos apresentados em conferências acadêmicas também podem ser relevantes para a fundamentação teórica, principalmente em áreas de pesquisa emergentes. Esses artigos são revisados por pares e muitas vezes discutem temas inovadores. O Código~\ref{cod:bibtex-inproceedings} apresenta o formato de citação BibTeX de um artigo de conferência.

% \begin{codigo}[htb]
% \legenda[cod:bibtex-inproceedings]{Estrutura BibTeX para citação de um artigo de conferência}
% \begin{lstlisting}
% @inproceedings{autor,
%     author    = {Sobrenome, Nome},
%     title     = {Título do Artigo},
%     booktitle = {Anais da Conferência},
%     year      = {Ano},
%     address   = {Local da Conferência},
%     pages     = {Páginas},
%     volume    = {Volume},
%     number    = {Número}
% }
% \end{lstlisting}
% \fonteautor
% \end{codigo}

% \subsection{Sites}

% Sites também podem ser utilizados como fontes em pesquisas acadêmicas, desde que sejam confiáveis e possuam relevância no contexto da pesquisa. O uso de fontes \textit{online} é comum em pesquisas que abordam temas emergentes ou questões práticas. A citação de sites deve incluir a data de acesso, como mostrado no Código~\ref{cod:bibtex-online}, uma vez que conteúdos online podem ser modificados.

% \begin{codigo}[htb]
% \legenda[cod:bibtex-online]{Estrutura BibTeX para citação de um site}
% \begin{lstlisting}
% @online{pagina,
%     author        = {Autor},
%     organization  = {Organização},
%     title         = {Título},
%     year          = {Ano},
%     url           = {URL da página},
%     urlaccessdate = {05 out. 2024}
% }
% \end{lstlisting}
% \fonteautor
% \end{codigo}

% A citação de verbetes da Wikipédia devem ser realizados conforme apresentado no Código~{cod:bibtex-wiki}. A citação será referenciada como apresentado em~\citeonline{wikipedia2024}.

% \begin{codigo}[htb]
% \legenda[cod:bibtex-wiki]{Estrutura BibTeX para citação de um artigo na Wikipédia}
% \begin{lstlisting}
% @inbook{wikipedia,
%     title          = {Machine Learning},
%     publisher      = {Wikimedia},
%     booktitle      = {Wikip\'edia: a enciclop\'edia livre},
%     year           = {2024},
%     url            = {https://en.wikipedia.org/wiki/Machine_learning},
%     urlaccessdate  = {05 out. 2024}
% }
% \end{lstlisting}
% \fonteautor
% \end{codigo}

% \section{Considerações Finais}

% A fundamentação teórica é a espinha dorsal de qualquer trabalho acadêmico, uma vez que é por meio dela que o autor constrói o suporte necessário para desenvolver suas ideias e hipóteses. O uso correto das citações e a inclusão de fontes confiáveis e pertinentes são fundamentais para assegurar a qualidade e a credibilidade da pesquisa.




% % -------------------------------------------------------------------------------------------------
% % -------------------------------------------------------------------------------------------------
% % -------------------------------------------------------------------------------------------------



% % -------------------------------------------------------------------------------------------------
% % -------------------------------------------------------------------------------------------------
% % 4. ILUSTRAÇÕES
% % -------------------------------------------------------------------------------------------------
% % -------------------------------------------------------------------------------------------------
% \chapter{Ilustrações e Gráficos}
% \label{cap:ilustracoes}

% Este capítulo aborda o uso adequado de figuras e subfiguras em trabalhos acadêmicos, além de orientações sobre a inserção e formatação de gráficos. O uso de elementos visuais, como figuras e gráficos, é essencial para a clareza da apresentação de dados e conceitos. Em pesquisas que demandam a comunicação visual de resultados, esses elementos complementam a argumentação textual e auxiliam na compreensão do conteúdo~\cite{tufte1990envisioning}.

% \section{Figuras}
% \label{sec:figuras}

% A inserção de figuras é feita utilizando o ambiente \texttt{figure}, que, em conjunto com o comando \texttt{\textbackslash fig}, permite a incorporação de imagens no documento. Cada figura deve ser acompanhada de uma legenda explicativa e, preferencialmente, ser referenciada no corpo do texto para que o leitor possa localizá-la facilmente~\cite{knuth1984texbook}. O Código~\ref{cod:figura} apresenta o exemplo de inserção de uma figura.

% \begin{codigo}[htb]
% \legenda[cod:figura]{Código para inclusão de uma figura no texto}
% \begin{lstlisting}
% \begin{figure}[htb]
%     \legenda[fig:figura]{Legenda da figura}
%     \fig{width=0.7\textwidth}{imagens/exemplo.png}
%     \fonteautor
% \end{figure}
% \end{lstlisting}
% \fonteautor
% \end{codigo}

% A Figura~\ref{fig:exemplo-fig} mostra um exemplo de figura usando os comandos do  Código~\ref{cod:figura}. 

% \begin{figure}[htb]
%     \legenda[fig:exemplo-fig]{Exemplo de figura}
%     \fig{width=0.65\textwidth}{img/figura}
%     \fonte{\citeauthoronline{xkcd-conclusion}, \citeyear{xkcd-conclusion}, Disponível em: \url{https://xkcd.com/1403/}}
% \end{figure}


% No Código~\ref{cod:figura}, a imagem é centralizada e redimensionada para 70\% da largura do texto. A legenda da figura é definida pelo comando \texttt{\textbackslash legenda}, e o rótulo é definido pelo parâmetro opcional (definido pelos colchetes). O rótulo permite que a figura seja referenciada no texto, como em ``veja a Figura~\ref{fig:exemplo-fig}''. 

% O parâmetro ``[htb]'' indica as preferências de posicionamento da figura: ``h'' para o local exato do código, ``t'' para o topo da página, e ``b'' para a parte inferior~\cite{goossens1997latex}. 

% \section{Subfiguras}
% \label{sec:subfiguras}

% As subfiguras são recomendadas em situações onde é necessário inserir múltiplas imagens agrupadas. Elas permitem organizar várias imagens em uma única figura, cada uma com sua legenda individual. Os comandos \texttt{\textbackslash sfig} e \texttt{\textbackslash lfig} podem ser utilizados para facilitar esse processo. O exemplo de código para inserção de subfiguras é apresentado no Código~\ref{cod:subfiguras}.

% \begin{codigo}[htb]
% \legenda[cod:subfiguras]{Código para inclusão de três subfiguras no texto}
% \begin{lstlisting}
% \begin{figure}[htb]
%     \legenda[fig:subfigura]{Exemplo de varias subfiguras}
%     % Linha 1: Figuras simples com label
%     \sfig[sfig:fig1]{scale=0.3}{imagens/fig1}\hfil
%     \sfig[sfig:fig2]{scale=0.3}{imagens/fig2}
%     % Linha em branco move as próximas subfiguras para a linha seguinte

%     % Linha 2: Figuras simples sem label
%     \sfig{scale=0.3}{imagens/fig1}\hfil
%     \sfig{scale=0.3}{imagens/fig2}

%     % Linha 3: Figuras com legenda e com label
%     \lfig[lfig:fig1]{scale=0.3}{imagens/fig1}{legenda 1}\hfil
%     \lfig[lfig:fig2]{scale=0.3}{imagens/fig2}{legenda 2}

%     % Linha 3: Figuras com legenda e sem label
%     \lfig{scale=0.3}{imagens/fig1}{legenda 1}\hfil
%     \lfig{scale=0.3}{imagens/fig2}{legenda 2}
%     \fonteautor
% \end{figure}
% \end{lstlisting}
% \fonteautor
% \end{codigo}

% No Código~\ref{cod:subfiguras}, três subfiguras são organizadas em uma linha, com espaçamento entre elas proporcionado pelo comando \textit{\textbackslash hfill}. As subfiguras podem ser referenciadas individualmente, como ``veja a Subfigura~\ref{sfig:fig1}'', ou em conjunto, como na ``veja a Figura~\ref{fig:subfigura1}''. Esse recurso é útil quando se deseja comparar diferentes imagens de maneira simultânea~\cite{goossens1997latex}.

% A Figura~\ref{fig:subfigura1} apresenta um exemplo composição de subfiguras usando o comando \texttt{\textbackslash sfig}. Neste exemplo, todas as subfiguras foram organizadas em apenas uma linha.

% \begin{figure}[htb]
%     \legenda[fig:subfigura1]{Exemplo de subfiguras simples}
%     \sfig[sfig:fig1]{scale=0.2}{img/fig1}\hfil
%     \sfig{scale=0.2}{img/fig2}\hfil
%     \sfig{scale=0.2}{img/fig3}\hfil
%     \sfig{scale=0.2}{img/fig4}
%     \fonte{\citeauthoronline{nintendo}, \citeyear{nintendo}, Disponível em: \url{https://www.nintendo.com}}
% \end{figure}

% O comando \texttt{\textbackslash lfig} deve ser utilizado quando for necessária a inclusão de legendas explicativas para cada uma das subfiguras. Este comando recebe um terceiro parâmetro que é a legenda referente a subfigura. A Figura~\ref{fig:subfigura2} mostra um exemplo usando este comando. Ambos os comandos podem ser combinados em uma mesma composição dentro do ambiente \texttt{figure}. 

% \begin{figure}[htb]
%     \legenda[fig:subfigura2]{Exemplo de subfiguras com legendas individuais}
%     \lfig[sfig:mario]{scale=0.35}{img/fig1}{Mário}\hfil
%     \lfig[sfig:luigi]{scale=0.35}{img/fig2}{Mário Verde}\hfil
%     \lfig[sfig:peach]{scale=0.35}{img/fig3}{Princesa Peach}

%     \lfig{scale=0.35}{img/fig4}{Cogumelo Toad}\hfil
%     \lfig{scale=0.35}{img/fig5}{Rei Koopa}\hfil
%     \lfig{scale=0.35}{img/fig6}{Yoshi}
%     \fonte{\citeauthoronline{nintendo}, \citeyear{nintendo}, Disponível em: \url{https://www.nintendo.com}}
% \end{figure}

% \section{Gráficos}

% A inserção de gráficos segue os mesmos princípios da inserção de figuras. O ambiente \texttt{grafico} é utilizado para acomodar gráficos gerados a partir de softwares de plotagem de dados ou imagens externas~\cite{mittelbach2004latex}. O Código~\ref{cod:grafico} apresenta um exemplo de inserção de gráfico e o Gráfico~\ref{fig:grafico} mostra o resultado do código.

% \begin{codigo}[htb]
% \legenda[cod:grafico]{Código para inclusão de um gráfico no texto}
% \begin{lstlisting}
% \begin{grafico}[htb]
%     \legenda[fig:grafico]{Legenda do grafico}
%     \fig{width=0.7\textwidth}{imagens/exemplo.png}
%     \fonteautor
% \end{grafico}
% \end{lstlisting}
% \fonteautor
% \end{codigo}

% \begin{grafico}[htb]
%     \legenda[fig:grafico]{Exemplo de gráfico}
%     \fig{width=0.7\textwidth}{img/grafico}
%     \fonte{\citeauthor{xkcd-graph}, \citeyear{xkcd-graph}, Disponível em: \url{https://xkcd.com/1945/}}
% \end{grafico}

% Gráficos são recursos visuais valiosos para a representação de dados quantitativos e a comunicação clara de tendências ou padrões observados durante a pesquisa. Assim como nas figuras, é importante que os gráficos sejam acompanhados de legendas claras e referenciados no corpo do texto~\cite{tufte1990envisioning}.

% \section{Considerações Finais}

% O uso adequado de figuras, subfiguras e gráficos em trabalhos acadêmicos é fundamental para uma comunicação visual eficiente. Esses elementos não apenas ilustram os conceitos discutidos, mas também ajudam a sintetizar e evidenciar os resultados da pesquisa. Quando bem utilizados, eles aumentam a clareza e a coesão do trabalho, facilitando a compreensão pelo leitor~\cite{tufte1990envisioning, knuth1984texbook}.



% % -------------------------------------------------------------------------------------------------
% % -------------------------------------------------------------------------------------------------
% % -------------------------------------------------------------------------------------------------



% % -------------------------------------------------------------------------------------------------
% % -------------------------------------------------------------------------------------------------
% % 5. TABELAS E QUADROS
% % -------------------------------------------------------------------------------------------------
% % -------------------------------------------------------------------------------------------------
% \chapter{Tabelas e Quadros}
% \label{cap:tabelas}

% Este capítulo aborda as técnicas de criação de tabelas e quadros, com a distinção de seus usos e estruturas. A apresentação organizada de dados é fundamental em trabalhos acadêmicos, e tanto as tabelas quanto os quadros desempenham um papel essencial na comunicação clara e precisa das informações. Contudo, eles diferem na natureza e no tipo de dado que apresentam.

% \section{Diferença entre Tabelas e Quadros}

% Embora os termos ``tabela'' e ``quadro'' sejam muitas vezes usados de forma intercambiável, há distinções claras entre esses elementos no contexto acadêmico. Tabelas são utilizadas principalmente para exibir dados quantitativos de forma estruturada, facilitando comparações entre diferentes valores ou categorias. Por outro lado, quadros organizam informações qualitativas, como definições ou classificações de conceitos, sendo mais descritivos e menos focados em dados numéricos~\cite{abnt2018}.

% Assim, enquanto as tabelas proporcionam uma apresentação visual de dados comparativos ou numéricos, os quadros são apropriados para descrever ou categorizar informações qualitativas, o que evidencia a diferença fundamental entre ambos~\cite{chagas2016}.

% \section{Tabelas}

% As tabelas são criadas utilizando-se os ambientes \texttt{table} e \texttt{tabular}. O primeiro define o posicionamento e a legenda, enquanto o segundo organiza o conteúdo em linhas e colunas. O Código~\ref{cod:tabela} apresenta uma tabela com quatro colunas, com diferentes tipos de alinhamento: fixo com~3cm~(p), à esquerda~(l), centralizado~(c) e à direita~(d). A Tabela~\ref{tab:tabela} apresenta o resultado da tabela descrita no código.

% \begin{codigo}[htb]
% \legenda[cod:tabela]{Código para criação de uma tabela}
% \begin{lstlisting}
% \begin{table}[htb]
%     \centering
%     \legenda[tab:tabela]{Exemplo de uma tabela}
%     \begin{tabular}{p{3cm}lcr}
%         \toprule
%         \textbf{Nome} & \textbf{Idade} & \textbf{Prova} & \textbf{Trabalhos} \\
%         \midrule
%         \midrule
%         Joao    & 24 & 8,5 & 4,8 \\
%         Maria   & 20 & 9,0 & 7,1 \\
%         Pedro   & 22 & 4,8 & 9,9 \\
%         \bottomrule
%     \end{tabular}
%     \fonteautor
% \end{table}
% \end{lstlisting}
% \fonteautor
% \end{codigo}

% \begin{table}[htb]
%     \centering
%     \legenda[tab:tabela]{Exemplo de uma tabela}
%     \begin{tabular}{p{3cm}lcr}
%         \toprule
%         \textbf{Nome} & \textbf{Idade} & \textbf{Prova} & \textbf{Trabalhos} \\
%         \midrule
%         \midrule
%         Joao    & 24 & 8,5 & 4,8 \\
%         Maria   & 20 & 9,0 & 7,1 \\
%         Pedro   & 22 & 4,8 & 9,9 \\
%         \bottomrule
%     \end{tabular}
%     \fonteautor
% \end{table}



% \section{Quadros}

% Quadros, por sua vez, são mais adequados para a organização de informações qualitativas, sendo amplamente utilizados para categorizações ou descrições. Embora utilizem a mesma estrutura básica de uma tabela, sua função é distinta. O Código~\ref{cod:quadro} mostra um exemplo de quadro no qual são categorizados tipos de pesquisa, e o Quadro~\ref{tab:quadro} mostra o resultado deste código.

% \begin{codigo}[htb]
% \legenda[cod:quadro]{Código para criação de um quadro}
% \begin{lstlisting}
% \begin{quadro}[htb]
%     \centering
%     \legenda[tab:quadro]{Exemplo de Quadro com tipos de pesquisa}
%     \begin{tabular}{|l|p{9cm}|}
%         \hline
%         \textbf{Tipo de Pesquisa} & \textbf{Descrição} \\ 
%         \hline\hline
%         Pesquisa Exploratória & Investigação inicial sobre um fenômeno, sem hipóteses definidas \\ \hline
%         Pesquisa Descritiva   & Caracterização detalhada de um fenômeno, com hipóteses exploradas \\ \hline
%         Pesquisa Explicativa  & Busca identificar as causas de um fenômeno \\         \hline
%     \end{tabular}
%     % Fix para ajustar a distancia da fonte (apenas para quadros)
%     \vspace{0.2cm} 
%     \fonteautor
% \end{quadro}
% \end{lstlisting}
% \fonteautor
% \end{codigo}


% \begin{quadro}[htb]
%     \centering
%     \legenda[tab:quadro]{Exemplo de Quadro com tipos de pesquisa}
%     \begin{tabular}{|l|p{9cm}|}
%         \hline
%         \textbf{Tipo de Pesquisa} & \textbf{Descrição} \\ 
%         \hline\hline
%         Pesquisa Exploratória & Investigação inicial sobre um fenômeno, sem hipóteses definidas \\ \hline
%         Pesquisa Descritiva   & Caracterização detalhada de um fenômeno, com hipóteses exploradas \\ \hline
%         Pesquisa Explicativa  & Busca identificar as causas de um fenômeno \\ \hline
%     \end{tabular}
%     \vspace{0.2cm} % Fix para ajustar a distância da fonte
%     \fonteautor
% \end{quadro}

% O posicionamento de tabelas e quadros é controlado por meio de argumentos opcionais nos ambientes \texttt{table} e \texttt{quadro}. As opções mais comuns incluem:

% \begin{alineas}
%     \item \texttt{h} (\textit{here}): posiciona o objeto o mais próximo possível do ponto onde foi inserido no código;
%     \item \texttt{t} (\textit{top}): coloca o objeto no topo da página;
%     \item \texttt{b} (\textit{bottom}): coloca o objeto na parte inferior da página;
%     \item \texttt{p} (\textit{page of floats}): insere o objeto em uma página dedicada apenas a elementos flutuantes, como tabelas e figuras.
% \end{alineas}

% Por exemplo, o argumento \texttt{[htb]}, utilizado nos Códigos~\ref{cod:tabela}~e~\ref{cod:quadro}, permite posicionar a tabela ``aqui''~(\textit{here}), mas com a flexibilidade de movê-la para o topo~(\textit{top}) ou parte inferior~(\textit{bottom}) da página, se necessário.

% \section{Considerações Finais}

% A correta utilização de tabelas e quadros é essencial para a apresentação organizada de dados e informações em trabalhos acadêmicos. Tabelas são mais apropriadas para dados quantitativos, enquanto quadros são utilizados para informações qualitativas e descritivas. A aplicação adequada desses elementos contribui significativamente para a clareza e profissionalismo do trabalho, facilitando a leitura e interpretação dos dados apresentados~\cite{knuth1984texbook}.


% % -------------------------------------------------------------------------------------------------
% % -------------------------------------------------------------------------------------------------
% % -------------------------------------------------------------------------------------------------



% % -------------------------------------------------------------------------------------------------
% % -------------------------------------------------------------------------------------------------
% % 6. CÓDIGOS
% % -------------------------------------------------------------------------------------------------
% % -------------------------------------------------------------------------------------------------
% \chapter{Códigos e Algoritmos}
% \label{cap:codigos}

% Este capítulo apresenta o uso de algoritmos e trechos de código, abordando desde a inclusão básica até a personalização de estilos e numeração. A formatação clara e precisa de algoritmos e códigos é fundamental em áreas como ciência da computação, matemática e engenharias, onde a transparência dos métodos e a reprodutibilidade são essenciais para a validade dos resultados apresentados~\cite{knuth1984texbook}.


% \section{Inclusão de Códigos}

% Para incluir códigos de programação, o pacote mais comum é o \texttt{listings}. Esse pacote suporta várias linguagens de programação e permite a personalização do estilo de exibição, incluindo o realce de sintaxe, a numeração de linhas e a adição de comentários. O ambiente \texttt{codigo} formata a posição do código e localização da legenda e da fonte.

% O Código~\ref{cod:codigo} apresenta um exemplo de código para inclusão de códigos-fonte e o Código~\ref{cod:fatorial} mostra o código renderizado.
% Este código implementa uma função recursiva em Python para calcular o fatorial de um número. A sintaxe da linguagem Python é destacada automaticamente e o código é numerado.

% \begin{codigo}[htb]
% \legenda[cod:codigo]{Código para inclusão de códigos-fonte}
% \begin{lstlisting}
% \begin{codigo}[htb]
%     \legenda[cod:codigo]{Código para inclusão de códigos-fonte}
%     \begin{lstlisting}[language=Python]
%     def fatorial(n):
%         if n == 0:
%             return 1
%         else:
%             return n * fatorial(n - 1)
%     \end{lstlisting }
%     \fonteautor
% \end{codigo}
% \end{lstlisting}
% \fonteautor
% \end{codigo}

% \begin{codigo}[htb]
% \legenda[cod:fatorial]{Função recursiva para calcular o fatorial de um número}
% \begin{lstlisting}[language=Python]
% def fatorial(n):
%     if n == 0:
%         return 1
%     else:
%         return n * fatorial(n - 1)
% \end{lstlisting}
% \fonteautor
% \end{codigo}


% \section{Inclusão de Algoritmos}

% A formatação de algoritmos é realizada de maneira análoga à inclusão de códigos, porém, utiliza o ambiente \texttt{algoritmo}. O Algoritmo~\ref{alg:fatorial} mostra um exemplo de algoritmo para o cálculo do fatorial de maneira recursiva.

% \begin{algoritmo}
% \legenda[alg:fatorial]{Algoritmo para o calcular o fatorial de um número}
% \begin{lstlisting}[mathescape]
% ENTRADA: Número inteiro não negativo $n$
% SAÍDA Valor de $n!$

% SE $n = 0$
%     RETORNA $1$
% SENÃO
%     RETURNA $n \times \text{fatorial}(n-1)$
% FIM-SE
% \end{lstlisting}
% \fonteautor
% \end{algoritmo}


% \section{Considerações Finais}

% O uso de algoritmos e códigos é um recurso poderoso para a formatação de trabalhos acadêmicos e científicos, especialmente em áreas técnicas. A personalização oferecida pelo LaTeX permite a adaptação dos estilos de exibição conforme as necessidades específicas do autor ou da instituição.



% % -------------------------------------------------------------------------------------------------
% % -------------------------------------------------------------------------------------------------
% % -------------------------------------------------------------------------------------------------



% % -------------------------------------------------------------------------------------------------
% % 7. EQUAÇÕES E FÓRMULAS
% % -------------------------------------------------------------------------------------------------
% \chapter{Equações e Fórmulas}
% \label{cap:equacoes}

% Este capítulo apresenta como criar e formatar equações, usar comandos específicos, a numeração e referência de equações, além de explorar exemplos de equações e sistemas de equações. As equações e fórmulas matemáticas desempenham um papel essencial em diversos campos do conhecimento, especialmente nas áreas de ciências exatas, engenharias e tecnologia. Em trabalhos acadêmicos, é fundamental que a apresentação dessas expressões seja clara e bem estruturada, facilitando a comunicação precisa de conceitos matemáticos.

% \section{Equações Simples}

% Para a inserção de equações matemáticas, utiliza-se o ambiente \texttt{equation}. Por padrão, as equações são centralizadas e numeradas automaticamente, permitindo a referência posterior no corpo do texto. O Código~\ref{cod:eqEinstein} exemplifica a inclusão de uma equação.

% \begin{codigo}[htb]
% \legenda[cod:eqEinstein]{Código para inclusão de equações}
% \begin{lstlisting}
% \begin{equation}
%     E = mc^2
%     \label{eq:einstein}
% \end{equation}
% \end{lstlisting}
% \fonteautor
% \end{codigo}

% A Equação~\eqref{eq:einstein} exibe a equação de Einstein. A referência a essa equação no texto pode ser realizada com o comando \texttt{\textbackslash ref}, como apresentado na frase ``conforme mostrado na Equação~\ref{eq:einstein}''.

% \begin{equation}
%     E = mc^2
%     \label{eq:einstein}
% \end{equation}

% \section{Sistemas de Equações}

% Para formatar sistemas de equações, utiliza-se o ambiente \texttt{\textbackslash cases}, que insere chaves ao lado do sistema, facilitando a organização de condições ou equações inter-relacionadas. O Código~\ref{cod:sistemas} exemplifica a inclusão de um equações alinhadas.

% \begin{codigo}[htb]
% \legenda[cod:sistemas]{Código para inclusão de equações}
% \begin{lstlisting}
% \begin{equation}
%     f(x) =
%     \begin{cases}
%         x^2 & \text{se } x \geq 0 \\
%         -x^2 & \text{se } x < 0
%     \end{cases}
%     \label{eq:sistemas}
% \end{equation}
% \end{lstlisting}
% \fonteautor
% \end{codigo}

% A Equação~\eqref{eq:sistemas} mostra um exemplo de sistemas de equações. Neste exemplo, a função \( f(x) \) é definida de acordo com duas condições diferentes, dependendo do valor de \( x \). O uso do ambiente``cases'' formata o sistema de forma clara e visualmente organizada.


% \begin{equation}
%     f(x) =
%     \begin{cases}
%         x^2 & \text{se } x \geq 0 \\
%         -x^2 & \text{se } x < 0
%     \end{cases}
%     \label{eq:sistemas}
% \end{equation}


% \section{Símbolos e Notações Avançadas}

% O \LaTeX\ oferece uma vasta gama de símbolos e notações matemáticas, permitindo a criação de expressões complexas de forma elegante. Entre os exemplos mais comuns, incluem-se:

% \begin{alineas}
%     \item Frações: \verb|\frac{a}{b}| para \( \frac{a}{b} \)
%     \item Somatórios: \verb|\sum_{i=1}^{n} i^2| para \( \sum_{i=1}^{n} i^2 \)
%     \item Integrais: \verb|\int_0^1 x^2 \, dx| para \( \int_0^1 x^2 \, dx \)
%     \item Vetores: \verb|\vec{v}| para \( \vec{v} \)
% \end{alineas}

% Através desses comandos, é possível construir fórmulas matemáticas de alta complexidade de maneira eficiente e esteticamente agradável.

% \section{Considerações Finais}

% O LaTeX é uma ferramenta robusta para a formatação de equações e fórmulas matemáticas, oferecendo recursos avançados que garantem precisão e consistência na apresentação de conceitos matemáticos. Com ambientes específicos para equações, sistemas de equações e expressões \textit{inline}, o autor possuis grande flexibilidade para a elaboração de textos acadêmicos de alta qualidade, facilitando a leitura e a compreensão dos leitores.



% % -------------------------------------------------------------------------------------------------
% % -------------------------------------------------------------------------------------------------
% % -------------------------------------------------------------------------------------------------



% % -------------------------------------------------------------------------------------------------
% % 7. ALÍNEAS
% % -------------------------------------------------------------------------------------------------
% \chapter{Alíneas}
% \label{cap:alineas}

% Este capítulo apresenta o modo de uso de alíneas. Em documentos acadêmicos, o uso correto de alíneas é fundamental para garantir a clareza na apresentação de informações que não possuem um título próprio. A norma técnica da Associação Brasileira de Normas Técnicas (ABNT) orienta a estruturação adequada das alíneas de forma a facilitar a compreensão e a organização textual. O uso das alíneas segue diretrizes específicas que devem ser rigorosamente aplicadas, conforme descrito a seguir~\cite{abnt2012}.

% As alíneas são empregadas quando é necessário enumerar vários assuntos em uma seção sem título. Um exemplo de alíneas pode ser observado na classificação dos paradigmas de programação. Os principais paradigmas de programação são:

% \begin{alineas}
%     \item Programação imperativa;
%     \item Programação funcional;
%     \item Programação orientada a objetos;
%     \item Programação lógica:
%         \begin{subalineas}
%             \item baseada no uso de lógica formal para expressar programas;
%             \item foca em resolver problemas declarando relações e permitindo que o computador derive soluções;
%             \item exemplos comuns incluem linguagens como Prolog.
%         \end{subalineas}
% \end{alineas}

% Neste exemplo, observa-se o uso correto de alíneas para organizar a informação de maneira hierárquica e clara. O uso de subalíneas foi aplicado para detalhar a ``Programação lógica''. Este tipo de estruturação facilita a compreensão e o destaque de elementos dentro de uma lista de assuntos.



% % -------------------------------------------------------------------------------------------------
% % -------------------------------------------------------------------------------------------------
% % -------------------------------------------------------------------------------------------------



% % -------------------------------------------------------------------------------------------------
% % 8. CRONOGRAMA
% % -------------------------------------------------------------------------------------------------
% \chapter{Cronograma}
% \label{cap:cronograma}

% O cronograma é uma ferramenta essencial no planejamento de um Trabalho de Conclusão de Curso~(TCC). Ele representa a programação das atividades necessárias para a conclusão do projeto, facilitando a organização e o gerenciamento do tempo. Um cronograma bem elaborado é fundamental para garantir que todas as etapas do TCC sejam realizadas dentro do prazo, evitando sobrecargas de trabalho.

% O cronograma deve incluir as principais etapas do TCC, como a definição do tema, revisão bibliográfica, desenvolvimento da metodologia, coleta e análise de dados, redação do trabalho e entrega final. O Quadro~\ref{tab:cronograma} apresenta um exemplo de cronograma.

% \begin{quadro}[h]
%     \legenda[tab:cronograma]{Cronograma detalhado das atividades do TCC}
%     \begin{tabular}{|p{7cm}|c|c|c|c|c|}
%         \hline
%         \textbf{Atividade}                      & \textbf{Ago} & \textbf{Set} & \textbf{Out} & \textbf{Nov} & \textbf{Dez} \\
%         \hline\hline
%         Definição do Tema                      & \cellcolor{black}  &  &  &  &  \\\hline
%         Revisão Bibliográfica                  & \cellcolor{black}  & \cellcolor{black}  &  &  &  \\\hline
%         Desenvolvimento da Metodologia         &  & \cellcolor{black}  & \cellcolor{black}  &  &  \\\hline
%         Coleta de Dados                        &  &  & \cellcolor{black}  & \cellcolor{black}  &  \\\hline
%         Análise dos Dados                      &  &  &  & \cellcolor{black}  & \cellcolor{black}  \\\hline
%         Redação do TCC                        &  &  &  & \cellcolor{black}  & \cellcolor{black}  \\\hline
%         Revisão Final e Entrega                &  &  &  &  & \cellcolor{black}  \\
%         \hline
%     \end{tabular}
%     \vspace{0.2cm}
%     \fonteautor
% \end{quadro}




% % -------------------------------------------------------------------------------------------------
% % 9. CONCLUSÃO
% % -------------------------------------------------------------------------------------------------
% \chapter{Conclusão}
% \label{cap:conclusao}

% A conclusão é uma das partes mais relevantes de um Trabalho de Conclusão de Curso~(TCC), pois é nela que o autor sintetiza os principais achados da pesquisa e apresenta suas reflexões finais. O objetivo principal da conclusão é proporcionar uma visão clara e objetiva sobre os resultados obtidos, além de contextualizá-los dentro da problemática abordada no trabalho. Para isso, uma conclusão eficaz deve seguir algumas diretrizes, garantindo que todos os pontos cruciais sejam abordados.

% Primeiramente, a conclusão deve retomar os objetivos do trabalho. Ao revisitar as perguntas de pesquisa e os objetivos estabelecidos na introdução, o autor consegue avaliar se foram alcançados ou não. Este processo de revisão não apenas reforça a relevância do tema, mas também permite que o leitor entenda como as informações apresentadas ao longo do texto se conectam com o que foi inicialmente proposto. Segundo \citeonline{gil}, a conclusão deve ``revisitar os objetivos traçados, avaliando o quanto se conseguiu avançar na pesquisa''~\cite{gil}.

% Em segundo lugar, é fundamental que a conclusão contenha uma síntese dos principais resultados e achados da pesquisa. Essa síntese deve ser objetiva e direta, evitando a inclusão de novos dados ou informações não discutidas previamente no corpo do texto. A análise crítica dos resultados obtidos pode ser feita de maneira a apontar as implicações e contribuições que o trabalho traz para o campo de estudo em questão. Como afirma~\citeauthoronline{lakatos}, ``a conclusão é o momento de apresentar as implicações do estudo realizado, bem como suas contribuições para a área de conhecimento''~\cite{lakatos}.

% Por fim, a conclusão pode incluir sugestões para futuras pesquisas. Este aspecto é importante porque pode abrir novas possibilidades de investigação, além de indicar áreas que precisam de mais atenção ou aprofundamento. A inclusão de sugestões é uma forma de reconhecer que o conhecimento é um processo contínuo e que existem lacunas a serem preenchidas. Como destaca \citeonline{lakatos}, ``sugerir novos caminhos de pesquisa é uma maneira de contribuir para o desenvolvimento da área e para a construção do conhecimento''~\cite{lakatos}.

% Em suma, a conclusão deve retomar os objetivos, sintetizar os resultados e sugerir novas pesquisas. Essa estrutura não apenas ajuda o leitor a entender o trabalho de forma mais clara, mas também enfatiza a contribuição do autor para a área estudada.



% -------------------------------------------------------------------------------------------------
% -------------------------------------------------------------------------------------------------
% -------------------------------------------------------------------------------------------------



% =================================================================================================
% -------------------------------------------------------------------------------------------------
% - ELEMENTOS PÓS-TEXTUAIS                                                                            -
% -------------------------------------------------------------------------------------------------
% =================================================================================================

% Inicia a seção pós-textual
\postextual

% -------------------------------------------------------------------------------------------------
% REFERÊNCIAS BIBLIOGRÁFICAS
% -------------------------------------------------------------------------------------------------

% Carrega o arquivo referencias.bib
\bibliography{referencias}


% % -------------------------------------------------------------------------------------------------
% % APÊNDICES
% % -------------------------------------------------------------------------------------------------
% \apendice[ape:exemplo]{Exemplo de Apêndice}

% Os apêndices e anexos são elementos complementares em um Trabalho de Conclusão de Curso~(TCC) que têm como objetivo fornecer informações adicionais relevantes ao entendimento da pesquisa. Embora sejam frequentemente utilizados de forma intercambiável, eles possuem finalidades distintas e devem ser utilizados em contextos específicos. A correta utilização de apêndices e anexos contribui para a clareza e a organização do trabalho acadêmico.

% Os apêndices são utilizados para incluir informações que são de autoria do próprio autor do TCC. Isso pode incluir questionários, roteiros de entrevistas, tabelas com dados adicionais e qualquer outro material que seja essencial para a compreensão do trabalho, mas que não é possível inserir diretamente no corpo do texto. Segundo~\citeonline{severino}, ``os apêndices servem para apresentar materiais elaborados pelo autor que complementam a sua pesquisa''~\cite{severino}. Dessa forma, o uso de apêndices é fundamental quando se deseja apresentar materiais que, embora relevantes, não se encaixam na narrativa principal do trabalho.


% % -------------------------------------------------------------------------------------------------
% % ANEXOS
% % -------------------------------------------------------------------------------------------------
% \anexo[ane:exemplo]{Exemplo de Anexo}

% Os anexos são documentos que não são de autoria do autor do TCC, mas que são pertinentes à pesquisa. Eles podem incluir gráficos, documentos oficiais, normativas e textos que foram utilizados como referência ou que fornecem suporte à pesquisa. Segundo a Associação Brasileira de Normas Técnicas (ABNT), ``os anexos são documentos que, embora não sejam parte integrante do texto, são relevantes para a compreensão do tema tratado''~\cite{abnt2012}. O uso de anexos é indicado quando se deseja fornecer ao leitor informações complementares que não foram geradas pelo autor, mas que ajudam a esclarecer ou contextualizar a pesquisa.

% Ambos, apêndices e anexos, devem ser listados em uma página específica após a conclusão do trabalho, de forma a facilitar a consulta por parte do leitor. A correta distinção entre apêndices e anexos, bem como o seu uso apropriado, é essencial para a organização e a credibilidade do TCC. Ao fornecer informações adicionais de maneira clara e sistemática, o autor contribui para uma melhor compreensão do seu trabalho, enriquecendo a experiência do leitor.


\end{document}